\documentclass[11pt]{article}

\usepackage[a4paper,margin=1in]{geometry}
\usepackage{amsmath,amssymb,amsthm,mathtools}
\usepackage[hidelinks]{hyperref}
\hypersetup{
  pdfauthor={Bangyen Pham},
  pdftitle={Every Row Is a Top Row: Coefficient Order Statistics of Multiples of a Power},
  pdfsubject={Exact coefficient order statistics of polynomial multiples of (x-r)^L},
  pdfkeywords={polynomial multiples, coefficient mass, reduced height, linear programming duality, exponential sums},
  pdflang={en-US}
}

\newtheorem{theorem}{Theorem}[section]
\newtheorem{lemma}[theorem]{Lemma}
\newtheorem{corollary}[theorem]{Corollary}
\newtheorem{proposition}[theorem]{Proposition}
\newtheorem*{mainthm}{Main Theorem}
\newcommand{\R}{\mathbb R}
\newcommand{\N}{\mathbb N}
\DeclareMathOperator{\Tail}{Tail}
\DeclareMathOperator{\sgn}{sgn}

\title{Every Row Is a Top Row: Coefficient Order Statistics of Multiples of a Power}
\author{\shortstack{Bangyen Pham\\
\small Independent Researcher\\
\small\href{mailto:bangyenp@gmail.com}{\texttt{bangyenp@gmail.com}}\\
\small\href{https://orcid.org/0009-0006-9928-2088}{ORCID: 0009-0006-9928-2088}}}
\date{September 2026}

\begin{document}
\maketitle

\begin{abstract}
For a real polynomial $F$ let $b_k(F)$ be the $k$-th largest of the
magnitudes of its nonleading coefficients, counted with multiplicity.
For $r>1$ let $V_r(L,k)$ be the infimum of $b_k(F)$ over monic multiples
$F$ of $(x-r)^L$.  We prove that for $r\ge2$ every row is a top row:
$V_r(L,k)=V_r(L-k+1,1)$ for $1\le k\le L$, so exempting the $k-1$ largest
coefficients costs exactly $k-1$ roots.  Linear-programming duality turns
each row into a weighted $\ell^1$ problem for polynomials with prescribed
integer zeros, and convexity together with first-order optimality shows
that each exempted position costs at most one root.  For $1<r<2$ the rows
are not all top rows.  On each diagonal $L-k=\mathrm{const}$ they equal
the values of explicit prefix-constrained problems in all but finitely
many rows, and in every row when the diagonal's top row is at most one.
Every row $k\ge2$ is the minimum of the preceding row on its diagonal and
finitely many terms, each attained by a polynomial with integer zeros
whose optimality, for rational $r$, is a finite exact check.  Exact
certificates at $r=\frac54$ and $r=\frac43$ show that the prefix
description fails in general and can fail in the third row of a diagonal
while holding in the second.
\end{abstract}

\noindent\textbf{Keywords.} Polynomial multiples; coefficient order
statistics; reduced height; linear programming duality.

\noindent\textbf{2020 Mathematics Subject Classification.} Primary 11C08; Secondary 26C05, 90C05.

\section{Introduction}\label{sec:intro}

Let $F=\sum_jf_jx^j$ be a nonzero real polynomial of degree $D$, and write
$b_k(F)$ for the $k$-th entry of its nonleading coefficient magnitudes
$|f_j|$, $j<D$, listed in decreasing order with multiplicity.  The companion
paper~\cite{coefficient-mass} proves, for real
roots $2\le r_1\le\cdots\le r_L$, that every nonzero multiple $F$ of
$\prod_i(x-r_i)$ has
\[
  b_k(F)\ \ge\ |f_D|\prod_{i=k}^L(r_i-1)\qquad(1\le k\le L)
\]
\cite[Theorem~3.2]{coefficient-mass}, and that at the root $2$ every monic
multiple of $(x-2)^L$ has $b_k(F)\ge L-k+1$
\cite[Theorem~4.13]{coefficient-mass}.  We call the infimum of
$b_k(F)$ over monic multiples $F$ of $\prod_i(x-r_i)$ the row $k$; the bounds
above are lower bounds for it.  This paper determines the rows when all
roots are equal.

\paragraph{The rows at a single root.}
Fix a real $r>1$ and put
\[
  V_r(L,k)=\inf\{b_k(F):\ F\text{ a monic real multiple of }(x-r)^L\},
  \qquad \beta_r(n)=V_r(n,1),
\]
so $\beta_r(n)$ is the top row of $(x-r)^n$.  Exempting the $k-1$ largest
coefficients always costs at least $k-1$ of the roots:
$V_r(L,k)\le\beta_r(L-k+1)$ (Proposition~\ref{prop:rowvalue}).  Our main result is that equality holds
at every root $r\ge2$.

\begin{mainthm}
For every real $r\ge2$ and all $1\le k\le L$,
$V_r(L,k)=\beta_r(L-k+1)$ (Theorem~\ref{thm:allrowstwo}).
\end{mainthm}

At $r=2$ this sharpens $b_k\ge L-k+1$ to $b_k\ge\beta_2(L-k+1)$; for instance
$q(s)=(1-\frac s2)(1-\frac s5)$ is admissible for $\beta_2(3)$ and has
$\sum_{s\ge1}|q(s)|2^{-s}=\frac{11}{40}$, so
$\beta_2(3)\ge\frac{40}{11}>3$.

For $1<r<2$ we compare the rows with the prefix values $\nu_i(n)$: the
least weighted $\ell^1$ norm $\sum_{s\ge1}|q(s)|r^{-s}$ of a polynomial $q$
with $q(0)=1$ and $\deg q\le n+i-2$ that vanishes on $[1,i-1]$.  There the
rows are not all top rows; the last row is
$V_r(L,L)=(r-1)^L<\beta_r(1)$ for $L\ge2$ (Corollary~\ref{cor:lastrowr}).
The prefix values satisfy
\[
  \frac1{\nu_k(n)+E_k(r)}\ \le\ V_r(L,k)\ \le\ \min_{1\le i\le k}\frac1{\nu_i(n)},
  \qquad n=L-k+1,
\]
with an explicit loss $E_k(r)$ from the first $2k-2$ positions
(Theorem~\ref{thm:prefixrows}).  For $1<r<2$ and fixed $n$ (indeed
$n=o(k/\log k)$), $E_k(r)/\nu_k(n)\to0$ as $k\to\infty$
(Corollary~\ref{cor:belowtwo}).

The prefix identity, equality in the upper bound, holds for every $k$
whenever $\beta_r(n)\le1$ (Theorem~\ref{thm:onepolyrows}), and on every
diagonal in all but finitely many rows (Corollary~\ref{cor:longdiagfinite}).
More precisely, it holds whenever $k\ge\max(K,\alpha(n-1))$, where
$\alpha\ge1$ is any constant with $\rho(\alpha)<1$ for an explicit function
$\rho$ depending on $r$, and $K$ depends on $r$ and the chosen $\alpha$
(Theorem~\ref{thm:longdiag}, Proposition~\ref{prop:longdiagexplicit}).

The prefix identity fails in general: exact certificates give
$V_{5/4}(16,2)=1/\nu_3(14)<\min_{i\le2}1/\nu_i(15)$ and
$V_{4/3}(19,2)<\min_{i\le2}1/\nu_i(18)$
(Proposition~\ref{prop:secondrowexact}).  They rest on a finite reduction:
the second row $V_r(n+1,2)$ is the minimum of the top row $\beta_r(n)$ and
finitely many exempted-position terms (Theorem~\ref{thm:secondrow}).  The
identity can fail in the second row even when $\nu_2(n)>\nu_1(n)$: at
$r=\frac{177}{167}$, $V_r(51,2)<1/\nu_2(50)<1/\nu_1(50)$
(Proposition~\ref{prop:secondrowhyp}).  Every
row $k\ge2$ is likewise the minimum of $V_r(L-1,k-1)$ and finitely many
exempted-set terms (Theorem~\ref{thm:finiterows}), and exact certificates
give $V_{4/3}(16,3)<V_{4/3}(15,2)=\min_{i\le3}1/\nu_i(14)$
(Proposition~\ref{prop:thirdrowexact}), so the identity can fail in the
third row of a diagonal while holding in the second.  The proofs of these three
propositions are computer-assisted: they list the certifying zero sets and
reduce the claims to finitely many comparisons of explicit rationals, which
the accompanying tests carry out in exact arithmetic.

\paragraph{Method.}
Divisibility by $(x-r)^L$ is a set of linear conditions on the
coefficients, and duality for linear programs turns each row into a problem
about polynomials
$q$ of degree less than $L$ with $q(0)=1$: minimize
\[
  \Phi_r(q)=\sum_{s\ge1}|q(s)|\,r^{-s}
\]
subject to $q$ vanishing on the set $S$ of backward distances of the
exempted positions (Proposition~\ref{prop:rowvalue}).  Minimizers exist
and can be taken with integer zeros, $q(s)=\prod_{z\in Z}(1-s/z)$
(Lemma~\ref{lem:vertexopt}).  Two facts then settle $r\ge2$.
\begin{itemize}
\item A convexity argument, crossing, shows that one insertion at the
first gap of a minimizer handles the last exempted position, wherever it lies
(Theorem~\ref{thm:crossing}).
\item First-order optimality of the minimizer in a single direction bounds
the cost of that insertion by the factor $\max\{1,1/(r-1)\}$
(Lemma~\ref{lem:optins}).
\end{itemize}
Below $2$ this factor exceeds $1$, and the proofs use instead one
polynomial for all exempted sets (Theorem~\ref{thm:onepoly}) and the
supermultiplicativity of the prefix values (Lemma~\ref{lem:prefixshift}).

\paragraph{Relation to the literature.}
The rows $k\ge2$, which exempt the largest coefficients, are an
order-statistic variant, introduced in the companion
paper~\cite{coefficient-mass}, of the reduced height.  The top row
$\beta_r(n)=V_r(n,1)$ is the reduced height of $(x-r)^n$ in the sense of
Dubickas and Jankauskas~\cite{dj}, the infimum of the heights of its monic
multiples; the reduced length~\cite{schinzel} is the $\ell^1$ analogue.
Lower bounds for the coefficients of polynomials with a root of high
multiplicity go back to the work of Beaucoup, Borwein, Boyd and Pinner on
multiple roots of power series with bounded
coefficients~\cite{bbbp-multiple}, which bounds the multiplicity at a
given point in terms of the coefficient size.

\paragraph{Organization.}
Section~\ref{sec:rowvalues} sets up the row values and their dual
problems, extends the tail bound of~\cite{coefficient-mass} to arbitrary
nodes, and determines the last row.  Section~\ref{sec:prefix} brackets
every row by prefix values.  Section~\ref{sec:crossing} shows that one
insertion covers every later exempted position,
Section~\ref{sec:allrowstwo} proves that every row is a top row at
$r\ge2$, and Section~\ref{sec:belowtwo} treats $1<r<2$.

\paragraph{Conventions.}
For a finite set $Z$ of positive integers we write
$r_Z(s)=\prod_{z\in Z}(1-s/z)$, so $r_Z(0)=1$ and $\deg r_Z=|Z|$;
Table~\ref{tab:notation} lists this and the other notation used in more
than one place, with the section or result where it is defined, and
symbols used only inside one proof are defined there.  All logarithms are
natural, $\N_{>0}=\{1,2,\ldots\}$, and $[u,v]$ denotes the integers in that
range.  Results of the companion paper are cited by their numbers there.

\begin{table}[tp]
\centering
\small
\renewcommand{\arraystretch}{1.15}
\begin{tabular}{@{}l p{0.47\textwidth} l@{}}
\hline
Symbol & Meaning & Defined in\\
\hline
$b_k(F)$ & $k$-th largest nonleading coefficient magnitude of $F$
  & Section~\ref{sec:intro}\\
$V_r(L,k)$ & row $k$: $\inf b_k(F)$ over monic multiples $F$ of $(x-r)^L$
  & Section~\ref{sec:rowvalues}\\
$\beta_r(n)$ & top row $V_r(n,1)$ & Section~\ref{sec:rowvalues}\\
$n$ & $L-k+1$; $\beta_r(n)$ is the top row of the diagonal $L-k=n-1$
  & Section~\ref{sec:rowvalues}\\
$a$ & $1/(r-1)$ & Section~\ref{sec:rowvalues}\\
$r_Z(s)$ & $\prod_{z\in Z}(1-s/z)$ & Conventions\\
$\Phi_r(q)$, $\Phi_r(Z)$ & $\sum_{s\ge1}|q(s)|r^{-s}$, and $\Phi_r(r_Z)$
  & Section~\ref{sec:rowvalues}\\
$\mu_r(S,L)$ & $\inf\Phi_r(q)$ over real $q$, $\deg q<L$, $q(0)=1$, $q|_S=0$
  & Section~\ref{sec:rowvalues}\\
worst exempted set & $S$ with $|S|=k-1$ and $V_r(L,k)=1/\mu_r(S,L)$
  & Section~\ref{sec:rowvalues}\\
$\Tail(w)$, leading run & $\sum_{d\ge1}|w_d|$; largest $\ell$ with
  $[1,\ell]\subseteq Z$ & Section~\ref{sec:rowvalues}\\
$\nu_i(n)$ & prefix value $\mu_r([1,i-1],n+i-1)$ & Section~\ref{sec:prefix}\\
$m_i(n)$ & $\nu_i(n)/a^i$, a negative-binomial expectation
  & Section~\ref{sec:prefix}\\
$W_k$ & pointwise bound for $|r_S|$ over $|S|\le k-1$
  & Section~\ref{sec:prefix}\\
$\Omega_k(n)$ & $\inf_p\sum_sW_k(s)r^{-s}|p(s)|$, $\deg p<n$, $p(0)=1$
  & Section~\ref{sec:prefix}\\
$E_k(r)$ & loss from the first $2k-2$ positions & Section~\ref{sec:prefix}\\
prefix identity & $V_r(L,k)=\min_{i\le k}1/\nu_i(n)$
  & \eqref{eq:prefixid}\\
$\mathrm{ins}(Z,c)$ & $Z$ with $c$ inserted and the elements from $c$ on
  moved up by one & Section~\ref{sec:crossing}\\
$h_t(W)$ & least integer above $t$ not in $W$ & Section~\ref{sec:crossing}\\
$\Phi_{r,D}$, $\mu_{r,D}$ & $\Phi_r$ and $\mu_r$ with the sum cut at $s\le D$
  & Section~\ref{sec:allrowstwo}\\
$\Phi_{r,D}^{<c}$, $\Phi_{r,D}^{>c}$ & the parts of $\Phi_{r,D}$ with
  $s<c$ and $s>c$ & Section~\ref{sec:allrowstwo}\\
$\omega_i(s)$ & $\binom{s-1}{i-1}r^{-s}$ & Section~\ref{sec:onepoly}\\
$\mathcal Q_n$ & real $p$ with $\deg p\le n-1$ and $p(0)=1$
  & Section~\ref{sec:onepoly}\\
$A_i(p)$ & $\sum_s\omega_i(s)|p(s)|=\Phi_r(r_{[1,i-1]}p)$
  & Section~\ref{sec:onepoly}\\
$\Lambda_k(n)$ & $\inf_{p\in\mathcal Q_n}\max_{i\le k}A_i(p)$
  & Theorem~\ref{thm:onepoly}\\
$w_k$ & $\omega_k-\omega_{k-1}$ & Section~\ref{sec:longdiag}\\
$\tau_k(n)$ & $\inf_{p\in\mathcal Q_n}\sum_{s\ge2k-2}w_k(s)|p(s)|$
  & Section~\ref{sec:longdiag}\\
$N_k(r)$ & largest excess of an $\omega_i$, $i<k$, over $\omega_k$ on
  $[1,2k-3]$ & Section~\ref{sec:longdiag}\\
$\theta$ & $4(r-1)/r^2$ & Corollary~\ref{cor:belowtwo}\\
$M_k$ & a mode of $\omega_k$ & Lemma~\ref{lem:nbmass}\\
$\Xi(k)$, $\rho(\alpha)$ & explicit bounds for long diagonals
  & Proposition~\ref{prop:longdiagexplicit}\\
$\eta_\sigma(n)$ & $\mu_r(\{\sigma\},n+1)$, one exempted position
  & Section~\ref{sec:secondrow}\\
$\Psi_r(p)$, $R_p(\sigma)$ & $\sum_ss|p(s)|r^{-s}$;
  $\sum_{s>\sigma}(s-\sigma)|p(s)|r^{-s}$ & Section~\ref{sec:secondrow}\\
$e_y$, $G_y$ & the vertex test & Lemma~\ref{lem:vertexopt}\\
$p_*$, $\sigma_*$ & a minimizer for $\nu_1(n)$; its tail threshold
  & Theorem~\ref{thm:secondrow}\\
$Y_n$, $Y'_n$, $Z_S(n)$ & named zero sets of the exact certificates
  & Section~\ref{sec:secondrow}\\
$\phi_m(s,t)$, $\Delta_{q,m}(t)$ & weight and gap for $m$ far zeros
  & Section~\ref{sec:finiterows}\\
$t_m(q)$ & least integer $t\ge1$ with $\Delta_{q,m}(t)\le0$
  & Lemma~\ref{lem:farzeros}\\
$T(N)$, $\mathcal S$ & thresholds and the finite tree of exempted sets
  & Theorem~\ref{thm:finiterows}\\
\hline
\end{tabular}
\caption{Notation.}\label{tab:notation}
\end{table}

\section{Row values and the dual problem}\label{sec:rowvalues}

Throughout the paper $r>1$ is real, and for real polynomials $q$ and
finite $Z\subset\N_{>0}$ we put
\[
  a=\frac1{r-1},\qquad
  \Phi_r(q)=\sum_{s\ge1}|q(s)|\,r^{-s},\qquad
  \Phi_r(Z)=\Phi_r(r_Z),
\]
and for finite $S\subset\N_{>0}$ and integers $L>|S|$ we let
$\mu_r(S,L)$ be the infimum of $\Phi_r(q)$ over real $q$ with $\deg q\le L-1$,
$q(0)=1$ and $q|_S=0$.  So $\Phi_2$ is the sum $\Phi$ of
\cite[Section~4.4]{coefficient-mass} and $\Phi_r(\varnothing)=a$.  For
$1\le k\le L$ put $n=L-k+1$ and
\[
  V_r(L,k)=\inf\{b_k(F): F\text{ a monic real multiple of }(x-r)^L\},
\]
the exact value of the row, and recall $\beta_r(n)=V_r(n,1)$.  At
$r=2$, \cite[Theorem~4.13]{coefficient-mass} says $V_2(L,k)\ge n$.

\begin{proposition}[The row values]\label{prop:rowvalue}
Let $r>1$ and $1\le k\le L$.
\begin{enumerate}
\item[(a)] Let $F=\sum_{j\le D}f_jx^j$ be a monic multiple of $(x-r)^L$,
let $J\subseteq[0,D-1]$ be any set of $k-1$ indices with
$\min_{j\in J}|f_j|\ge\max\{|f_j|:j<D,\ j\notin J\}$ (so $J$ carries $k-1$
largest nonleading magnitudes, ties broken arbitrarily), and let
$S=\{D-j:j\in J\}$ be the set of their backward distances from the leading
position.  If $q$ is real with $\deg q\le L-1$, $q(0)=1$ and $q|_S=0$, then
$b_k(F)\,\Phi_r(q)\ge1$.  In particular $b_k(F)\ge1/\Phi_r(Z)$ for every
finite $Z\supseteq S$ with $|Z|\le L-1$.
\item[(b)] $V_r(L,k)=\inf_{|S|=k-1}1/\mu_r(S,L)$, the infimum over
$S\subset\N_{>0}$.
\item[(c)] $\beta_r(n)=1/\mu_r(\varnothing,n)$, and $V_r(L,k)\le\beta_r(n)$.
\end{enumerate}
\end{proposition}

\begin{proof}
(a) This is the proof of \cite[Lemma~4.5]{coefficient-mass} with $2$ replaced by $r$:
$(q(D-\delta)F)(r)=0$ with $\delta=x\,d/dx$ gives
$\sum_{j\le D}f_jq(D-j)r^{j-D}=0$, so
$1\le\sum_{s\ge1,\,s\notin S}|f_{D-s}||q(s)|r^{-s}\le b_k(F)\Phi_r(q)$.
The last step uses $|f_j|\le b_k(F)$ for $j<D$, $j\notin J$: the $k-1$
magnitudes on $J$ are at least every other nonleading one, so $b_k(F)$ is
the largest magnitude off $J$ (here $D\ge L\ge k$).  Only $|S|=k-1$ and
this inequality are used below, so every result stated for ``the set $S$
of (a)'' holds for every admissible choice of $J$.

(b) By (a), $b_k(F)\ge1/\mu_r(S_F,L)$ for the set $S_F$ of (a), which
gives ``$\ge$''.  For ``$\le$'' fix $S$ with $|S|=k-1$ and $D\ge\max(L,\max S)$.
Writing $F=\sum_{s=0}^Dg_sx^{D-s}$ with $g_0=1$, divisibility by
$(x-r)^L$ is the set of $L$ linear conditions
$\sum_{s=0}^Dg_sq(s)r^{-s}=0$ for $\deg q\le L-1$ (as in (a), these are the
conditions $(\delta^iF)(r)=0$, $i<L$).  Minimizing
$\max_{s\in[1,D]\setminus S}|g_s|$ over this affine set is a finite linear
program, feasible because $D\ge L$, and by linear-programming duality its value is
$1/\mu_{r,D}(S,L)$, where $\mu_{r,D}$ is defined like $\mu_r$ with the sum
$\Phi_r(q)$ truncated to $s\le D$ (the dual variables are the values
$q(s)r^{-s}$ of a polynomial $q$ of degree at most $L-1$ vanishing on $S$;
for weak duality, $|q(0)|=|\sum_{s\ge1,\,s\notin S}g_sq(s)r^{-s}|$).  Since
the $k-1$ positions at distances $S$ can carry at most the $k-1$ largest
magnitudes, $b_k(F)\le\max_{s\notin S}|g_s|$, so
$V_r(L,k)\le1/\mu_{r,D}(S,L)$ for every $D$.  Finally
$\mu_{r,D}(S,L)\to\mu_r(S,L)$: the truncated sums increase with $D$ and
are at most $\Phi_r$; on the affine space $\{\deg q\le L-1,\ q(0)=1\}$ the
values $q(1),\ldots,q(L-1)$ are coordinates, so minimizers $q_D$ of the
truncated problems ($D\ge L$) have bounded coefficients, and a limit point
$q_*$ satisfies $q_*(0)=1$, $q_*|_S=0$ and, for each fixed $M$,
$\sum_{s\le M}|q_*(s)|r^{-s}\le\lim_D\mu_{r,D}(S,L)$; letting $M\to\infty$
gives $\mu_r(S,L)\le\Phi_r(q_*)\le\lim_D\mu_{r,D}(S,L)$.

(c) For $k=1$, (b) with $S=\varnothing$ is $\beta_r(n)=V_r(n,1)=1/\mu_r(\varnothing,n)$.
Now let $S_N=[N+1,N+k-1]$.  A polynomial $q$ of degree at most $L-1$ with
$q(0)=1$ and $q|_{S_N}=0$ factors as $q=q_0\pi_N$ with
$\pi_N(s)=\prod_{\sigma\in S_N}(1-s/\sigma)$, $\deg q_0\le n-1$ and
$q_0(0)=1$.  For $1\le s\le M<N$ we have $\pi_N(s)\ge(1-M/N)^{k-1}$, so
$\Phi_r(q)\ge(1-M/N)^{k-1}\mu_{r,M}(\varnothing,n)$.  Hence
$\liminf_N\mu_r(S_N,L)\ge\mu_{r,M}(\varnothing,n)$ for every $M$, and by
the convergence in (b) it is at least $\mu_r(\varnothing,n)$.  By (b),
$V_r(L,k)\le\inf_N1/\mu_r(S_N,L)\le1/\mu_r(\varnothing,n)=\beta_r(n)$.
\end{proof}

So every row is at most the top row of a multiple of $(x-r)^n$, the value
approached when the $k-1$ exempted positions lie far below the leading one.
We call $S\subset\N_{>0}$ with $|S|=k-1$ a \emph{worst exempted set} for
$V_r(L,k)$ if it attains the infimum in Proposition~\ref{prop:rowvalue}(b), that is,
$V_r(L,k)=1/\mu_r(S,L)$, and for $k=2$ we call its element a \emph{worst
exempted position}.  Such a set need not exist or be unique.

We next extend the displaced-zero tail bound
\cite[Theorem~2.2]{coefficient-mass}, proved there for nodes
$y_i\le\frac12$, to all nodes in $(0,1)$ (Theorem~\ref{thm:tailgen}), and
at the root $r$ we deduce that deleting the largest zero costs at most the
factor $\max(1,a)$ (Lemma~\ref{lem:confdelr}).  Such certificates give
lower bounds; Lemma~\ref{lem:confdelr} is used in the remark after
Theorem~\ref{thm:crossing}, while the last row
(Corollary~\ref{cor:lastrowr}) needs only convexity.

As in \cite[Section~2]{coefficient-mass}, for distinct nodes
$y_1,\ldots,y_c\in(0,1)$ and $Z\subset\N_{>0}$ with $|Z|=c-1$, the
\emph{exponential sum} $w$ on these nodes with $w_0=1$ and $w|_Z=0$ is the
sequence $w_d=\sum_i\gamma_iy_i^d$ with these values, $\Tail(w)=\sum_{d\ge1}|w_d|$,
and the \emph{leading run} of $Z$ is the largest $\ell\ge0$ with
$[1,\ell]\subseteq Z$.

\begin{theorem}[Tail bound at arbitrary nodes]\label{thm:tailgen}
Let $0<y_1<\cdots<y_c<1$, let $Z\subset\N_{>0}$ with $|Z|=c-1$ have leading
run $\ell$, let $w$ be the exponential sum on these nodes with $w_0=1$ and
$w|_Z=0$, and put $\varphi_i=y_i/(1-y_i)$.  Then
\[
  \Tail(w)\le\prod_{i=1}^{\ell+1}\varphi_i\prod_{i=\ell+2}^c\max(1,\varphi_i).
\]
Moreover, if $Z\ne[1,c-1]$ and $u$ is the sum on $y_1,\ldots,y_{c-1}$ with
$u_0=1$ and $u|_{Z\setminus\{\max Z\}}=0$, then
$\Tail(w)\le\max(1,\varphi_c)\Tail(u)$.
\end{theorem}

\begin{proof}
We use the zero bound of \cite[Section~2]{coefficient-mass}: a nonzero
exponential sum on $n$ distinct positive nodes has at most $n-1$ real
zeros, counted with multiplicity, so a nonzero sum on these nodes with
$n-1$ prescribed distinct real zeros has exactly those zeros, each simple
and a change of sign.  In particular a kernel vector of the $c$ conditions
defining $w$ would be a nonzero $c$-node sum with the $c$ zeros
$\{0\}\cup Z$, so $w$ exists and is unique; this needs only distinct
positive nodes.  As in \cite[Section~2]{coefficient-mass}, sums are
regarded as functions of the real variable $d\ge0$.  We induct on
$c-1-\ell$.

\emph{Base case.}  If $Z=[1,c-1]$, then $\ell=c-1$ and
\cite[Lemma~2.1]{coefficient-mass}, which needs no cutoff, gives
$\Tail(w)=\prod_{i=1}^c\varphi_i$.

\emph{Displaced step.}  Let $\ell<c-1$.  Then $Z\ne[1,c-1]$, so $c\ge2$ and
$z=\max Z\ge c\ge\ell+2$.  Put $Z'=Z\setminus\{z\}$ and
$m=\max(Z'\cup\{0\})<z$, and let $u$ be as in the statement:
the sum on $y_1,\ldots,y_{c-1}$ with $u_0=1$ and $u|_{Z'}=0$.  Let $G$ be a
nonzero sum on all $c$ nodes vanishing on $\{0\}\cup Z'$, which exists
since these are $c-1$ linear conditions on $c$ coefficients; $G_z\ne0$, as
otherwise $G$ would have $c$ zeros, so we normalize $G_z=1$.  Put
$\xi=|u_z|$.  If $\xi=0$, then $u$, a sum on the $c$ nodes with coefficient
$0$ at $y_c$, satisfies the conditions defining $w$, so $w=u$ and
$\Tail(w)\le\max(1,\varphi_c)\Tail(u)$.  Otherwise let $\eta=\sgn(u_z)$; then $w=u-\eta\xi G$,
both sides satisfying the same $c$ conditions.

By the zero bound, $G$ has exactly the simple zeros $\{0\}\cup Z'$, $u$
exactly $Z'$, and $w$ exactly $Z$.  So $u/G$ has constant sign on
$(0,\infty)\setminus Z'$, equal to its sign $\eta$ at $z$, and
$w_d=\eta\sgn(G_d)\psi(d)$ there, where $\psi(d)=|u_d|-\xi|G_d|$.  The
function $\psi$ is continuous, vanishes on $Z'$, has $\psi(0)=1$ and has no
zero in $(0,\min Z)$, where $w$ has none; so $\psi>0$ there.  Off $Z$, the
sign of $\psi$ does not change across a point of $Z'$, where $w$ and $G$
both change sign, and changes across $z$, where only $w$ does.  Hence
$|u_d|>\xi|G_d|$ for $1\le d<z$, $d\notin Z'$, and $|u_d|<\xi|G_d|$ for
$d>z$, with $|w_d|=\bigl||u_d|-\xi|G_d|\bigr|$ for all $d\ge1$.  Splitting
the tails at $z$ and using $|u_z|=\xi|G_z|$ gives
\[
  \Tail(u)-\Tail(w)=\xi\sum_{1\le d<z}|G_d|
  +\sum_{d\ge z}\bigl(2|u_d|-\xi|G_d|\bigr),
\]
all series converging because the nodes are below $1$.  Therefore
\begin{equation}\label{eq:tailsplit}
  \Tail(w)\le\Tail(u)+\xi\sum_{d\ge z}|G_d|-2\sum_{d\ge z}|u_d|.
\end{equation}

We bound the middle term.  Put $\hat u=\eta u$, $\tau=G_z/\hat u_z=1/\xi$,
and $E_d=G_d-y_cG_{d-1}$ for $d\ge1$.  By the zero bound, $\hat u$ and $G$
are positive for $d>m$: neither has a zero there, and $\hat u_z,G_z>0$.
Writing $G_d=\sum_ig_iy_i^d$ gives
$E_d=\sum_{i<c}g_i(1-y_c/y_i)y_i^d$, a sum on $y_1,\ldots,y_{c-1}$.  With
the marked points $\{0\}\cup Z'$ and $\widehat V=\hat u$,
\cite[Lemma~2.3]{coefficient-mass}, which needs no upper bound on the
nodes, gives $E_d\le\sigma\hat u_d$ for $d\ge m+1$, where
$\sigma=G_{m+1}/\hat u_{m+1}$.  We claim $\sigma\le\tau$.  If $z=m+1$ they
are equal.  If $z>m+1$, the $c$-node sum $\Theta=G-\tau\hat u$ vanishes on
the $c-1$ points $Z'\cup\{z\}$ and has $\Theta_0=-\tau\eta\ne0$, so it
changes sign exactly at those points.  The $c-2$ points of $Z'$ lie in
$(0,m+1)$, so $\sgn\Theta_{m+1}=(-1)^c\sgn\Theta_0=-(-1)^c\eta$, and
$\eta=(-1)^c$ because $u_0=1$ and $u$ changes sign at the $c-2$ points of
$Z'$, all below $z$.  So $\Theta_{m+1}<0$, that is, $\sigma<\tau$.  Hence
$E_d\le\tau\hat u_d$ for $d\ge z$, and induction on $t$ in
$G_{z+t}=y_cG_{z+t-1}+E_{z+t}$, starting from $G_z=\tau\hat u_z$, gives
\[
  G_{z+t}\le\tau\sum_{j=0}^ty_c^{\,t-j}\hat u_{z+j}\qquad(t\ge0).
\]
Summing over $t$,
$\xi\sum_{d\ge z}|G_d|\le\frac1{1-y_c}\sum_{d\ge z}|u_d|$, and
\eqref{eq:tailsplit} becomes
\[
  \Tail(w)\le\Tail(u)+\Bigl(\frac1{1-y_c}-2\Bigr)\sum_{d\ge z}|u_d|.
\]
If $y_c\le\frac12$ the bracket is at most $0$ and $\Tail(w)\le\Tail(u)$.
If $y_c>\frac12$ the bracket is $\varphi_c-1>0$, and
$\sum_{d\ge z}|u_d|\le\Tail(u)$ gives $\Tail(w)\le\varphi_c\Tail(u)$.  In
both cases $\Tail(w)\le\max(1,\varphi_c)\Tail(u)$.  Finally $z>\ell+1$, so
$Z'$ has the same leading run $\ell$ as $Z$, and $\ell\le c-2=|Z'|$; the
induction hypothesis for $u$ on $y_1,\ldots,y_{c-1}$ bounds $\Tail(u)$ by
$\prod_{i=1}^{\ell+1}\varphi_i\prod_{i=\ell+2}^{c-1}\max(1,\varphi_i)$,
which gives the bound for $w$.
\end{proof}

For $y_c\le\frac12$ every $\max(1,\varphi_i)$ is $1$ and this is
\cite[Theorem~2.2]{coefficient-mass}.

\begin{lemma}[Deleting the largest zero at the root $r$]\label{lem:confdelr}
Let $r>1$.  Then $\Phi_r([1,\ell])=a^{\ell+1}$ for $\ell\ge0$, and for every
nonempty finite $Z$,
\[
  \Phi_r(Z)\le\max(1,a)\,\Phi_r\bigl(Z\setminus\{\max Z\}\bigr).
\]
Consequently $\Phi_r(Z)\le\Phi_r(Z\cap[1,t])$ for all $t$ when $r\ge2$, and
$\Phi_r(Z)\le a^{|Z|+1}$ when $1<r\le2$.
\end{lemma}

\begin{proof}
For $x=1/r$ we have $\sum_{s>\ell}\binom{s-1}\ell x^s=(x/(1-x))^{\ell+1}$
and $x/(1-x)=a$; since $|r_{[1,\ell]}(s)|=\binom{s-1}\ell$ for $s>\ell$ and
$r_{[1,\ell]}$ vanishes on $[1,\ell]$, this is $\Phi_r([1,\ell])=a^{\ell+1}$.
It settles the inequality for $Z=[1,|Z|]$, the ratio being $a$.

Now let $Z\ne[1,|Z|]$, put $c=|Z|+1$ and $Z'=Z\setminus\{\max Z\}$.  For
small $\varepsilon>0$ take the nodes $y_i=\frac1r-(c-i)\varepsilon$, so
$0<y_1<\cdots<y_c=\frac1r<1$ and $\varphi_c=a$.  Let $w^\varepsilon$ be the
sum on $y_1,\ldots,y_c$ with $w_0=1$ and $w|_Z=0$, and $u^\varepsilon$ the
sum on $y_1,\ldots,y_{c-1}$ with $u_0=1$ and $u|_{Z'}=0$.  Since
$Z\ne[1,c-1]$, Theorem~\ref{thm:tailgen} gives
$\Tail(w^\varepsilon)\le\max(1,a)\Tail(u^\varepsilon)$.

Let $\varepsilon\to0$.  The divided difference of $y\mapsto y^d$ at
$y_1,\ldots,y_i$ is $N_i(d)=h_{d-i+1}(y_1,\ldots,y_i)$, with $h_j$ the
complete homogeneous symmetric polynomial ($h_j=0$ for $j<0$), and it is a
combination of $y_1^d,\ldots,y_i^d$ with nonzero coefficient on $y_i^d$.  So
$w^\varepsilon_d=\sum_{i=1}^c\alpha_iN_i(d)$ for unique $\alpha_i$, and the
$c$ conditions at $d\in\{0\}\cup Z$ form a linear system for $\alpha$ whose
matrix $(N_i(d))$ tends to the matrix with entries
$\binom d{i-1}r^{-(d-i+1)}$, the value of $h_{d-i+1}$ at $i$ arguments equal
to $\frac1r$.  The limit matrix is invertible: after scaling rows by $r^d$
and columns by $r^{-(i-1)}$ its entries are $\binom d{i-1}$, and the
polynomials $\binom s{i-1}$, $i\le c$, form a basis of the polynomials of
degree below $c$, none of which but $0$ vanishes at the $c$ points
$\{0\}\cup Z$.  Hence $\alpha$ converges, to $\alpha^0$ say.  Since the $y_j$
are positive and at most $\frac1r$,
$0\le h_{d-i+1}(y_1,\ldots,y_i)\le\binom d{i-1}r^{-(d-i+1)}$, so
$|w^\varepsilon_d|\le C(d+1)^{c-1}r^{-d}$ with $C$ independent of
$\varepsilon$ and $d$.  The limit is $w^0_d=p(d)r^{-d}$ with
$p(s)=\sum_i\alpha^0_ir^{i-1}\binom s{i-1}$, so $\deg p\le|Z|$, $p(0)=1$
and $p|_Z=0$, which forces $p=r_Z$.  By dominated convergence
$\Tail(w^\varepsilon)\to\Phi_r(Z)$, and in the same way, on the nodes
$y_1,\ldots,y_{c-1}$, $\Tail(u^\varepsilon)\to\Phi_r(Z')$.  This proves the
inequality.

The consequences follow by deleting the largest element repeatedly: for
$r\ge2$ each step does not increase $\Phi_r$, and for $r\le2$ each of the
$|Z|$ steps from $Z$ to $\varnothing$ costs at most the factor $a$, with
$\Phi_r(\varnothing)=a$.
\end{proof}

At $r=2$, where $a=1$, this is \cite[Lemma~4.6]{coefficient-mass}.

\begin{corollary}[The last row]\label{cor:lastrowr}
$V_r(L,L)=r-1$ for $r\ge2$, and $V_r(L,L)=(r-1)^L$ for $1<r\le2$.
\end{corollary}

\begin{proof}
Lower bounds: Proposition~\ref{prop:rowvalue}(a) with $Z=S$, $|S|=L-1$, and
Lemma~\ref{lem:prefixpush} of Section~\ref{sec:onepoly} with $p=1$, whose
proof uses only the convexity of Lemma~\ref{lem:crossing}(b).  It gives
$\Phi_r(S)\le\max_{1\le i\le L}\Phi_r([1,i-1])=\max_{1\le i\le L}a^i$, by
the first claim of Lemma~\ref{lem:confdelr} (a direct computation), and
this maximum is $a$ for $r\ge2$ and $a^L$ for $r\le2$.  Upper bounds: for $r\ge2$,
Proposition~\ref{prop:rowvalue}(c) with $n=1$ and
$\beta_r(1)=1/\Phi_r(\varnothing)=r-1$; for $r\le2$, Proposition~\ref{prop:rowvalue}(b)
with $S=[1,L-1]$, where $q=r_S$ is forced and $\mu_r(S,L)=a^L$.
\end{proof}

At $r=2$ both values are $1$.

\section{Prefix values}\label{sec:prefix}

This section keeps the notation of the previous one and brackets every row
by prefix values (Theorem~\ref{thm:prefixrows}).  For $1<r<2$ the bracket
identifies the rows up to a factor tending to $1$ once $k$ is large
compared with $n$ (Corollary~\ref{cor:belowtwo}).

For $i\ge1$ and $n\ge1$ put
\[
  \nu_i(n)=\inf\Bigl\{\sum_{s\ge1}\binom{s-1}{i-1}r^{-s}|p(s)|:\
  p\in\R[s],\ \deg p\le n-1,\ p(0)=1\Bigr\},
\]
so $\nu_1(n)=\mu_r(\varnothing,n)=1/\beta_r(n)$ (Proposition~\ref{prop:rowvalue}(c)).  Since
$|r_{[1,i-1]}(s)|=\binom{s-1}{i-1}$ for $s\ge1$, we have
$\nu_i(n)=\mu_r([1,i-1],i+n-1)$.

The normalized prefix values $m_i(n)=\nu_i(n)/a^i$ are negative-binomial
expectations: $m_i(n)=\inf_p\mathbb E|p(X_i)|$, where $X_i$ counts the
trials up to the $i$-th success at success probability $1-1/r$, because
$\binom{s-1}{i-1}r^{-s}=a^i\Pr(X_i=s)$.  So $1/\nu_i(n)$ is a top-row value
for a transformed measure.

The weight
\[
  W_k(s)=\binom{s-1}{\min(k-1,\lfloor(s-1)/2\rfloor)}
  =\max\Bigl\{\binom{s-1}j:0\le j\le\min(k-1,s-1)\Bigr\},
\]
the second form by unimodality of the binomial coefficients, bounds
$|r_S(s)|$ for every $S$ with $|S|\le k-1$ (Lemma~\ref{lem:rowpointwise}).

The weighted problem
$\Omega_k(n)=\inf_p\sum_{s\ge1}W_k(s)r^{-s}|p(s)|$, with $p$ as in $\nu_i$,
gives the lower bound for the rows.

The loss
\[
  E_k(r)=\sum_{s=1}^{2k-2}\binom{s-1}{\lfloor(s-1)/2\rfloor}r^{-s}
\]
collects the first $2k-2$ positions, where $W_k$ can exceed
$\binom{s-1}{k-1}$.

\begin{lemma}[Integer zero sets suffice]\label{lem:intzeros}
Let $i,n\ge1$ and let $\omega(s)\ge0$ ($s\ge1$) vanish for $s<i$, be
positive for $s\ge i$, and satisfy $\sum_s\omega(s)s^{n-1}<\infty$.  Then
$\inf_p\sum_s\omega(s)|p(s)|$, over real $p$ with $\deg p\le n-1$ and
$p(0)=1$, equals the infimum over the polynomials
$p=\prod_{y\in Y}(1-s/y)$ with $Y$ a set of $n-1$ integers at least $i$.
\end{lemma}

\begin{proof}
Let $D\ge i+n$ and $f_D(p)=\sum_{s=i}^D\omega(s)|p(s)|$ on the affine space
$\mathcal A=\{\deg p\le n-1,\ p(0)=1\}$ of dimension $n-1$.  On each cell
$\{p\in\mathcal A:\ \epsilon_sp(s)\ge0,\ i\le s\le D\}$
($\epsilon_s=\pm1$), $f_D$ is linear and nonnegative.  Each cell is
pointed, because a polynomial of degree at most $n-1$ that vanishes at $0$
and at the $D-i+1\ge n$ points $[i,D]$ is zero.  So the minimum of $f_D$ over
a cell is attained at a vertex.  A vertex satisfies $n-1$ linearly
independent equations $p(y)=0$ with $y\in[i,D]$, so
$p=\prod_{y\in Y}(1-s/y)$ with $|Y|=n-1$.

Hence some such $p_D$ has $f_D(p_D)=\min_{\mathcal A}f_D\le
\sum_s\omega(s)|p(s)|$ for all $p\in\mathcal A$.  Since
$|p_D(s)|\le(1+s)^{n-1}$, the full sum at $p_D$ exceeds $f_D(p_D)$ by at
most $\sum_{s>D}\omega(s)(1+s)^{n-1}$, which tends to $0$.  So the
infimum lies between $\min f_D$ and $\min f_D$ plus this tail, and the
truncated minima increase
to it.
\end{proof}

\begin{lemma}[A uniform bound on the exempted factor]\label{lem:rowpointwise}
Let $S\subset\N_{>0}$ be finite, $s\ge1$, and $j_s=|S\cap[1,s-1]|$.  Then
$|r_S(s)|\le\binom{s-1}{j_s}$.  If $|S|\le k-1$, then
$|r_S(s)|\le W_k(s)$, and $W_k(s)=\binom{s-1}{k-1}$ for $s\ge2k-1$.
\end{lemma}

\begin{proof}
Let $\sigma_1<\cdots<\sigma_{j}$ be the elements of $S$ below $s$, so
$\sigma_l\ge l$ and $j=j_s\le s-1$.  The factor $(s-x)/x$ decreases in $x$, so
$|1-s/\sigma_l|\le(s-l)/l$.  Every $\sigma\in S$ with $\sigma\ge s$ has
$0\le1-s/\sigma<1$.  Hence
$|r_S(s)|\le\prod_{l\le j}(s-l)/l=\binom{s-1}j$.  The rest follows from
$j_s\le\min(k-1,s-1)$, unimodality, and the fact that
$\lfloor(s-1)/2\rfloor\ge k-1$ exactly when $s\ge2k-1$.
\end{proof}

\begin{theorem}[Rows through prefix sets]\label{thm:prefixrows}
Let $r>1$, $1\le k\le L$ and $n=L-k+1$.  Then
\begin{equation}\label{eq:prefixsandwich}
  \frac1{\nu_k(n)+E_k(r)}\ \le\ \frac1{\Omega_k(n)}\ \le\ V_r(L,k)\ \le\
  \min_{1\le i\le k}\frac1{\nu_i(n)} .
\end{equation}
\end{theorem}

\begin{proof}
\emph{Upper bound.}  Fix $1\le i\le k$ and, for large $N$, put
$S_N=[1,i-1]\cup[N+1,N+k-i]$, so $|S_N|=k-1$.  A real $q$ with
$\deg q\le L-1$, $q(0)=1$ and $q|_{S_N}=0$ factors as
$q=r_{[1,i-1]}\pi_Nq_0$ with $\pi_N(s)=\prod_{\sigma=N+1}^{N+k-i}(1-s/\sigma)$,
$\deg q_0\le n-1$ and $q_0(0)=1$.  For $s\le M<N$ we have
$\pi_N(s)\ge(1-M/N)^{k-i}$, so
\[
  \Phi_r(q)\ge\Bigl(1-\frac MN\Bigr)^{k-i}\nu_{i,M}(n),
\]
where $\nu_{i,M}$ is $\nu_i$ with the sum cut at $s\le M$.
Proposition~\ref{prop:rowvalue}(b) gives
$V_r(L,k)\le1/\mu_r(S_N,L)$.  Let $N\to\infty$, then $M\to\infty$.  The
proof of Lemma~\ref{lem:intzeros} gives $\nu_{i,M}(n)\to\nu_i(n)$, so
$V_r(L,k)\le1/\nu_i(n)$.

\emph{Lower bound.}  Let $F$ and $S$ be as in
Proposition~\ref{prop:rowvalue}(a), so $|S|=k-1$.  For any $p$ as in
$\Omega_k$, the polynomial $q=r_Sp$ is admissible there, and by
Lemma~\ref{lem:rowpointwise}
\[
  \Phi_r(q)\le\sum_sW_k(s)r^{-s}|p(s)| .
\]
Taking the infimum over $p$ gives $b_k(F)\ge1/\Omega_k(n)$.

Now take $p=\prod_{y\in Y}(1-s/y)$ with $|Y|=n-1$ and $Y\subset[k,\infty)$.
For $1\le s\le2k-2\le2y$ every factor satisfies $|1-s/y|\le1$, so
$|p(s)|\le1$.  For $s\ge2k-1$, $W_k(s)=\binom{s-1}{k-1}$.  Hence
\[
  \sum_sW_k(s)r^{-s}|p(s)|\le E_k(r)+\sum_{s\ge1}\binom{s-1}{k-1}r^{-s}|p(s)| .
\]
By Lemma~\ref{lem:intzeros} with $\omega(s)=\binom{s-1}{k-1}r^{-s}$ and
$i=k$, the infimum of the last sum over such $p$ is $\nu_k(n)$.  This gives
$\Omega_k(n)\le\nu_k(n)+E_k(r)$.
\end{proof}

The upper bound says that the prefix sets $S=[1,i-1]$, completed by a far
block, are admissible exempted sets.  The lower bound compares every $S$
with the prefix $[1,k-1]$.  The loss $E_k(r)$ comes from the first $2k-2$
positions only.  For $r<2$ and $n$ small compared with $k$ it is exponentially small
against $\nu_k(n)$ (Corollary~\ref{cor:belowtwo}).

\begin{lemma}[The negative-binomial mass]\label{lem:nbmass}
Let $r>1$, $k,n\ge1$, $\varpi=1-1/r$, and let $M_k$ be a mode of the
negative-binomial law $P(t)=\binom{t-1}{k-1}\varpi^k r^{-(t-k)}$, $t\ge k$.
Then $M_k\le(k-1)r/(r-1)+1$, and
\[
  \frac{\nu_k(n)}{a^k}=m_k(n)\ \ge\
  \frac{(n-1)!\,P(M_k)}{r^{n-1}\bigl(2(M_k+n-1)\bigr)^{n-1}},
  \qquad
  P(M_k)\ge\frac3{4\bigl(4\sqrt{kr}/(r-1)+1\bigr)} .
\]
\end{lemma}

\begin{proof}
First, $\binom{s-1}{k-1}r^{-s}=a^kP(s)$, which gives $\nu_k=a^km_k$.  The
ratio $P(t+1)/P(t)=t/(r(t-k+1))$ is at least $1/r$.  It is at least $1$
exactly when $t\le(k-1)/\varpi$, which locates the mode.

Let $t_l=M_k+l$ for $0\le l\le n-1$, and let $\ell_l$ be the Lagrange basis
on these nodes.  Then
$|\ell_l(0)|=\prod_{l'\ne l}t_{l'}/|l-l'|\le(M_k+n-1)^{n-1}/(l!\,(n-1-l)!)$.
Evaluating $1=p(0)=\sum_l\ell_l(0)p(t_l)$ gives
$\max_l|p(t_l)|\ge(n-1)!/(2(M_k+n-1))^{n-1}$.  Moreover
$P(t_l)\ge r^{-(n-1)}P(M_k)$.

The law has variance $kr/(r-1)^2$.  By Chebyshev's inequality at
least $\frac34$ of its mass lies within $2\sqrt{kr}/(r-1)$ of the mean, on at most
$4\sqrt{kr}/(r-1)+1$ integers.  So the mode carries at least $3/(4(4\sqrt{kr}/(r-1)+1))$.
\end{proof}

\begin{corollary}[Rows for roots between one and two]\label{cor:belowtwo}
Let $1<r<2$ and $\theta=4(r-1)/r^2<1$.  For all $1\le k\le L$ put
$n=L-k+1$ and $\varepsilon=E_k(r)/\nu_k(n)$.  Then
\[
  1\ \ge\ \frac{V_r(L,k)\,m_k(n)}{(r-1)^k}\ \ge\
  \frac1{1+\varepsilon},\qquad
  \varepsilon\le\frac{r^2\,\theta^k}{4(2-r)\,m_k(n)} .
\]
For $n\le k$, Lemma~\ref{lem:nbmass} and $M_k+n-1\le k(2r-1)/(r-1)$ give
\[
  \varepsilon\le\frac{r^2\bigl(4\sqrt r/(r-1)+1\bigr)}{3(2-r)}\,
  \frac{\bigl(2r(2r-1)k/(r-1)\bigr)^{n-1}}{(n-1)!}\,\sqrt k\,\theta^k .
\]
So for fixed $n$ (indeed for $n=o(k/\log k)$),
$V_r(L,k)=(1+o(1))(r-1)^k/m_k(n)=(1+o(1))/\mu_r([1,k-1],L)$ as
$k\to\infty$; in this sense the prefix $S=[1,k-1]$ is asymptotically a
worst exempted set.
\end{corollary}

\begin{proof}
Theorem~\ref{thm:prefixrows} with $i=k$ gives the upper bound, and
$1/(\nu_k+E_k)=(1/\nu_k)/(1+\varepsilon)$ gives the lower one, as
$1/\nu_k(n)=(r-1)^k/m_k(n)$.  Since
$\binom{s-1}{\lfloor(s-1)/2\rfloor}\le2^{s-1}$ and $2/r>1$,
\[
  E_k(r)\le\frac12\sum_{s=1}^{2k-2}\Bigl(\frac2r\Bigr)^s\le
  \frac{(2/r)^{2k-1}}{2(2/r-1)}=\frac{r^2}{4(2-r)}\Bigl(\frac4{r^2}\Bigr)^k .
\]
Finally $\varepsilon=E_k/\nu_k=(r-1)^kE_k/m_k$, which gives the bound on
$\varepsilon$; the last form uses $\nu_k(n)=\mu_r([1,k-1],L)$.  And
$\theta<1$ is $(r-2)^2>0$.
\end{proof}

We say that the \emph{prefix identity} holds at $(r,L,k)$ if the upper
bound of~\eqref{eq:prefixsandwich} is attained, that is, if
$V_r(L,k)$ equals $\min_{i\le k}1/\nu_i(n)$ with $n=L-k+1$:
\begin{equation}\label{eq:prefixid}
  V_r(L,k)=\min_{1\le i\le k}\frac1{\nu_i(n)}\qquad(n=L-k+1).
\end{equation}
For $r\ge2$ the loss $E_k(r)$ is not small, and the lower bound
of~\eqref{eq:prefixsandwich} does not give the identity.  It is false for some $1<r<2$ (Proposition~\ref{prop:secondrowexact}).
For $r\ge2$ it follows from Theorem~\ref{thm:allrowstwo} (see the remark
after it).  For $1<r<2$, Theorems~\ref{thm:onepolyrows}
and~\ref{thm:longdiag} prove it in two regimes, in both with the minimum at
$i=k$.

\section{Crossing an exempted position}\label{sec:crossing}

This section shows that the last exempted position costs at most one
insertion, made at a place that does not depend on that position.  It
rests on two facts: the polynomials admissible for
Proposition~\ref{prop:rowvalue}(a) with $\Phi_r\le\varphi$ form a convex
set, so a position is covered as soon as two admissible polynomials take
opposite signs there (Lemma~\ref{lem:crossing}); and an insertion reverses
the sign of $r_W$ at every later position that lies in neither set
(Lemma~\ref{lem:insflip}).

For finite $Z\subset\N_{>0}$ and an integer $c\ge1$ let
\[
  \mathrm{ins}(Z,c)=\{z\in Z:z<c\}\cup\{c\}\cup\{z+1:z\in Z,\ z\ge c\}.
\]

For a finite $W\subset\N_{>0}$ and an integer $t\ge0$ put
\[
  h_t(W)=\min\{x\in\N:\ x>t,\ x\notin W\},
\]
the first gap of $W$ above $t$.  Every integer in $(t,h_t(W))$ lies in $W$.
If $h>\max W$ then $\mathrm{ins}(W,h)=W\cup\{h\}$.

\begin{lemma}[Crossing]\label{lem:crossing}
Let $r>1$, $L\ge1$, $\varphi>0$, let $T\subset\N_{>0}$ be finite and
$\sigma\in\N_{>0}\setminus T$.
\begin{enumerate}
\item[(a)] Let $q_1,q_2$ be real polynomials with $\deg q_i\le L-1$,
$q_i(0)=1$, $q_i|_T=0$ and $\Phi_r(q_i)\le\varphi$.  If
$q_1(\sigma)q_2(\sigma)\le0$, then $|T|+1<L$ and
$\mu_r(T\cup\{\sigma\},L)\le\varphi$.
\item[(b)] For every real polynomial $Q$ and all reals $0<y\le\sigma'$,
\[
  \Phi_r\bigl(Q\cdot(1-s/\sigma')\bigr)\le
  \max\bigl\{\Phi_r(Q),\ \Phi_r\bigl(Q\cdot(1-s/y)\bigr)\bigr\}.
\]
\end{enumerate}
\end{lemma}

\begin{proof}
(a) If $q_1(\sigma)=0$ put $q=q_1$, and if $q_2(\sigma)=0$ put $q=q_2$.
Otherwise $\lambda=q_2(\sigma)/(q_2(\sigma)-q_1(\sigma))$ lies in $(0,1)$;
put $q=\lambda q_1+(1-\lambda)q_2$.  Then $q(\sigma)=0$, $q(0)=1$, $q|_T=0$
and $\deg q\le L-1$.  Moreover
$|q(s)|\le\lambda|q_1(s)|+(1-\lambda)|q_2(s)|$ for every $s$, so
$\Phi_r(q)\le\varphi$.  The polynomial $q$ is nonzero and vanishes at the
$|T|+1$ points of $T\cup\{\sigma\}$, so $|T|+1\le L-1$, and $q$ is admissible
for $\mu_r(T\cup\{\sigma\},L)$.

(b) With $\lambda=y/\sigma'\in(0,1]$ we have
$1-s/\sigma'=(1-\lambda)+\lambda(1-s/y)$.  So
$|Q(s)(1-s/\sigma')|\le(1-\lambda)|Q(s)|+\lambda|Q(s)(1-s/y)|$ for every $s$.
Multiply by $r^{-s}$ and sum.
\end{proof}

\begin{lemma}[An insertion reverses the later signs]\label{lem:insflip}
Let $W\subset\N_{>0}$ be finite, let $h\in\N_{>0}\setminus W$ and
$W'=\mathrm{ins}(W,h)$.  Then:
\begin{itemize}
\item $W'\cap[1,h-1]=W\cap[1,h-1]$, $h\in W'$ and $|W'|=|W|+1$;
\item for every integer $\sigma>h$ with $\sigma\notin W\cup W'$,
$r_W(\sigma)\,r_{W'}(\sigma)<0$.
\end{itemize}
\end{lemma}

\begin{proof}
The first three claims are the definition of $\mathrm{ins}$ (as $h\notin W$,
the elements of $W$ that are at least $h$ exceed $h$).  For a finite $Z$ and
an integer $\sigma\notin Z$, the factor $1-\sigma/z$ of $r_Z(\sigma)$ is
negative exactly when $z<\sigma$, so $r_Z(\sigma)$ has the sign
$(-1)^{|Z\cap[1,\sigma-1]|}$.

Let $\sigma>h$ with $\sigma\notin W\cup W'$, and put
$A=|W\cap[1,h-1]|$ and $B=|W\cap[h+1,\sigma-2]|$.  The elements of $W'$
above $h$ are the $w+1$ with $w\in W$, $w>h$.  Those below $\sigma$ come
from $w\le\sigma-2$, so $|W'\cap[1,\sigma-1]|=A+1+B$.  Next,
$|W\cap[1,\sigma-1]|=A+B+[\sigma-1\in W]$, since $h\notin W$.  If
$\sigma-1\in W$, then $\sigma-1\ne h$, so $\sigma-1>h$ and $\sigma\in W'$,
which is excluded.  So $|W\cap[1,\sigma-1]|=A+B$, and the two counts differ
by one.
\end{proof}

\begin{theorem}[One insertion covers every later position]\label{thm:crossing}
Let $r>1$, let $t\ge0$ be an integer, let $W\subset\N_{>0}$ be finite, let
$T\subseteq W\cap[1,t]$ and $h=h_t(W)$.  Then for every integer $\sigma>t$,
\[
  \mu_r\bigl(T\cup\{\sigma\},\,|W|+2\bigr)\ \le\
  \max\bigl\{\Phi_r(W),\ \Phi_r(\mathrm{ins}(W,h))\bigr\}.
\]
Consequently, let $F$ be a monic multiple of $(x-r)^L$, let $S$ be the
exempted set of Proposition~\ref{prop:rowvalue}(a) with $|S|=k-1\ge1$, and
let $t=\max(S\setminus\{\max S\})$ ($t=0$ if $k=2$).  If $W\supseteq S\setminus\{\max S\}$
and $|W|\le L-2$, then
\[
  b_k(F)\ \ge\ \frac1{\max\{\Phi_r(W),\,\Phi_r(\mathrm{ins}(W,h_t(W)))\}} .
\]
\end{theorem}

\begin{proof}
Put $W'=\mathrm{ins}(W,h)$.  The polynomials $r_W$ and $r_{W'}$ have degree
at most $|W|+1$, take the value $1$ at $0$, and vanish on $T$, because
$T\subseteq W\cap[1,t]\subseteq W\cap[1,h-1]\subseteq W'$.  If $\sigma<h$,
then $\sigma\in(t,h)\subseteq W$ and $r_W$ vanishes at $\sigma$.  If
$\sigma=h$, then $r_{W'}$ vanishes at $\sigma$.  If $\sigma>h$ and
$\sigma\in W\cup W'$, one of the two vanishes at $\sigma$.  Otherwise
$r_W(\sigma)r_{W'}(\sigma)<0$ by Lemma~\ref{lem:insflip}.  In every case
Lemma~\ref{lem:crossing}(a), with $L=|W|+2$ and
$\varphi=\max\{\Phi_r(W),\Phi_r(W')\}$, gives the first claim.

For the second, apply the first with $T=S\setminus\{\max S\}$ and
$\sigma=\max S$.  This gives polynomials $q$ with $\deg q\le|W|+1\le L-1$,
$q(0)=1$, $q|_S=0$ and $\Phi_r(q)$ arbitrarily close to
$\mu_r(S,|W|+2)$.  Proposition~\ref{prop:rowvalue}(a) gives
$b_k(F)\Phi_r(q)\ge1$ for each of them.
\end{proof}

The theorem, which needs no condition on $r$, reduces the last exempted
position to one insertion at the first gap above the previous ones, and
Section~\ref{sec:allrowstwo} bounds its cost when $W$ is a minimizer.
When $h>\max W$ the second set is $W\cup\{h\}$, and for $r\ge2$
Lemma~\ref{lem:confdelr} gives $\Phi_r(W\cup\{h\})\le\Phi_r(W)$.  So a set
$W$ without gaps in $(t,\max W]$ handles every later position for free.

\section{Every row is a top row at roots at least \texorpdfstring{$2$}{2}}\label{sec:allrowstwo}

This section proves the prefix identity~\eqref{eq:prefixid} at every root $r\ge2$
(Theorem~\ref{thm:allrowstwo}).  By Theorem~\ref{thm:crossing} it suffices
to bound the cost of one insertion at the first gap of a minimizer, and we
do this, in a truncated form, by a first-order optimality argument.

An insertion at a gap $c$ of a minimizer $W$ is paid for by the first
moment $\sum_ss|r_W(s)|r^{-s}$, which the optimality of $W$ bounds.  The
zeros of $W$ above $c$ move up by one, and above $c$ the new polynomial is
essentially the old one shifted by one position.  Measured against the
first moment, the shift costs at most the factor $1/r$.  The first moment
is in turn bounded by optimality, in the direction that trades the leading
coefficient of $r_{\mathrm{ins}(W,c)}$ against that of $s\,r_W$.

For $D\in\N_{>0}\cup\{\infty\}$ and a real polynomial $q$ put
\[
  \Phi_{r,D}(q)=\sum_{s=1}^{D}|q(s)|\,r^{-s},\qquad
  \Phi_{r,D}^{<c}(q)=\sum_{1\le s<c,\ s\le D}|q(s)|\,r^{-s},\qquad
  \Phi_{r,D}^{>c}(q)=\sum_{c<s\le D}|q(s)|\,r^{-s},
\]
and let $\mu_{r,D}(S,L)$ be the infimum of $\Phi_{r,D}(q)$ over real $q$ with
$\deg q\le L-1$, $q(0)=1$ and $q|_S=0$.  This is the truncated quantity of the
proof of Proposition~\ref{prop:rowvalue}(b), and
$\Phi_{r,\infty}=\Phi_r$, $\mu_{r,\infty}=\mu_r$.  For a finite
$W\subset\N_{>0}$ and an integer $c\ge1$ with $c\notin W$ put
$A_c=r_{W\cap[1,c-1]}$ and $B_c=r_{W\cap(c,\infty)}$, so $r_W=A_cB_c$.  Since
$\mathrm{ins}(W,c)=(W\cap[1,c-1])\cup\{c\}\cup\{w+1:w\in W,\ w>c\}$ and
$1-s/(b+1)=(1-(s-1)/b)\,b/(b+1)$,
\[
  r_{\mathrm{ins}(W,c)}(s)=A_c(s)\Bigl(1-\frac sc\Bigr)\frac{B_c(s-1)}{B_c(-1)},
  \qquad B_c(-1)=\prod_{w\in W,\ w>c}\Bigl(1+\frac1w\Bigr).
\]

\begin{lemma}[Vertex minimizers of truncated problems]\label{lem:truncvertex}
Let $r>1$, let $S\subset\N_{>0}$ be finite, let $L>|S|$, and let $D$ be an
integer with $D\ge L+\max(S\cup\{0\})$.  Then some $W$ with
$S\subseteq W\subseteq[1,D]$ and $|W|=L-1$ satisfies
$\Phi_{r,D}(r_W)=\mu_{r,D}(S,L)$.
\end{lemma}

\begin{proof}
This is the argument of Lemma~\ref{lem:intzeros}.  Put $d=L-1-|S|$ and
$N=[1,D]\setminus S$, so $|N|\ge D-|S|\ge L>d$.  Every admissible $q$
factors as $q=r_Sp$ with $\deg p\le d$ and $p(0)=1$.  Moreover
$\Phi_{r,D}(q)=f(p):=\sum_{s\in N}\omega(s)|p(s)|$ with
$\omega(s)=|r_S(s)|r^{-s}>0$ on $N$.  If $d=0$, then $p=1$ and $W=S$.  Let
$d\ge1$.  For each sign vector $(\epsilon_s)_{s\in N}$, the cell
$\{p:\deg p\le d,\ p(0)=1,\ \epsilon_sp(s)\ge0\ (s\in N)\}$ contains no line.
Indeed, a direction $e$ of a line has $e(0)=0$, $\deg e\le d$ and $e|_N=0$,
so $e=0$.  On the cell $f$ is affine and nonnegative, so its minimum there is
attained at a vertex.  A vertex satisfies $d$ linearly independent equations
$p(y)=0$ with $y\in N$.  So $p=\prod_{y\in Y}(1-s/y)$ for a set $Y\subseteq N$
of $d$ integers.  The finitely many cells cover all admissible $p$.  So the
minimum of $f$ is attained at such a $p$, and $W=S\cup Y$ works.
\end{proof}

\begin{lemma}[Inserting at a gap of a minimizer]\label{lem:optins}
Let $r>1$ and $D\in\N_{>0}\cup\{\infty\}$.  Let $T\subseteq W$ be finite
subsets of $\N_{>0}$, put $L=|W|+1$, and suppose that
$\varphi:=\Phi_{r,D}(r_W)=\mu_{r,D}(T,L)$.  Let $c\ge1$ be an integer with
$c\notin W$ and $c>\max T$ ($\max\varnothing=0$), and put
$Q=r_{\mathrm{ins}(W,c)}$.  Then
\[
  \Phi_{r,D}(Q)\le\frac{\varphi+(r-2)\,\Phi_{r,D}^{<c}(Q)}{r-1},
  \qquad
  \Phi_{r,D}^{<c}(Q)\le\Phi_{r,D}^{<c}(r_W)\le\varphi .
\]
In particular $\Phi_{r,D}(Q)\le\max\{1,a\}\,\varphi$, and
$\Phi_{r,D}(Q)\le\varphi$ if $r\ge2$.
\end{lemma}

\begin{proof}
Write $P=r_W$, $A=A_c$, $B=B_c$ and $\kappa=1/(cB(-1))>0$.  The zeros of $A$
lie in $[1,c-1]$, and those of $B$ are integers greater than $c$.  So $B>0$ on
$[-1,c)$, and $B(s-1)B(s)\ge0$ for every integer $s$, since $B$ does not
change sign between consecutive integers.

\emph{Step 1: signs.}  Let $1\le s<c$.  Then $P(s)Q(s)=A(s)^2(1-s/c)B(s)B(s-1)/B(-1)\ge0$.  If
$A(s)\ne0$, then
\[
  \frac{|Q(s)|}{|P(s)|}=\frac{c-s}c\prod_{w\in W,\ w>c}g(w),\qquad
  g(w)=\frac{(w+1-s)\,w}{(w-s)(w+1)}=1+\frac s{(w-s)(w+1)} .
\]
The factor $g$ exceeds $1$ and decreases in $w$.  The $m$ elements of $W$
above $c$ are at least $c+1,\ldots,c+m$, so the product is at most
$\prod_{j=1}^mg(c+j)=\frac{(c+m+1-s)(c+1)}{(c+1-s)(c+m+1)}\le\frac{c+1}{c+1-s}$.
Since $(c-s)(c+1)\le c(c+1-s)$, we get $|Q(s)|\le|P(s)|$.  If $A(s)=0$, then
$P(s)=Q(s)=0$.  So $Q$ has the sign of $P$ below $c$, and
$\Phi_{r,D}^{<c}(Q)\le\Phi_{r,D}^{<c}(P)\le\varphi$.  Next, $Q(c)=0$.  For
$s>c$ we have $P(s)Q(s)=A(s)^2(1-s/c)B(s)B(s-1)/B(-1)\le0$.

\emph{Step 2: the shift.}  Let $s>c$ and $u=s-1\ge c$.  Every factor
$(u+1-z)/(u-z)$ with $z\le c-1$ exceeds $1$, so
\[
  \frac{|A(u+1)|}{|A(u)|}=\prod_{z\in W\cap[1,c-1]}\frac{u+1-z}{u-z}
  \le\prod_{z=1}^{c-1}\frac{u+1-z}{u-z}=\frac u{u+1-c} .
\]
Hence $|A(s)|(s-c)\le(s-1)|A(s-1)|$, and
$|Q(s)|r^{-s}\le\frac\kappa r(s-1)|P(s-1)|r^{-(s-1)}$.  Put
$M=\sum_{1\le u\le D}u|P(u)|r^{-u}$, which is finite, and split it as
$M=M_<+M_\ge$ according to $u<c$ or $u\ge c$.  Summing over $c<s\le D$ gives
\[
  \Phi_{r,D}^{>c}(Q)\le\frac\kappa r\,M_\ge .
\]

\emph{Step 3: optimality.}  The leading coefficient of $Q$ is $-\kappa$ times
that of $P$, and $\deg Q=\deg P+1=L$.  So $H=P-Q-\kappa sP$ has
$\deg H\le L-1$ and $H(0)=0$.  It vanishes on $T$, because $T\subseteq W\cap[1,c-1]$
consists of zeros of $A$.  Hence, for $\varepsilon>0$, the polynomial
\[
  q_\varepsilon=P+\varepsilon H=(1+\varepsilon-\varepsilon\kappa s)P-\varepsilon Q
\]
is admissible for $\mu_{r,D}(T,L)$, and $\Phi_{r,D}(q_\varepsilon)\ge\varphi$.
Let $\varepsilon\le1/(\kappa c)$.  For $1\le s<c$ we have
$1+\varepsilon-\varepsilon\kappa s\ge\varepsilon$, and by Step~1
$|q_\varepsilon(s)|=(1+\varepsilon-\varepsilon\kappa s)|P(s)|-\varepsilon|Q(s)|$.
For $s\ge c$, $|q_\varepsilon(s)|\le|1+\varepsilon-\varepsilon\kappa s||P(s)|+\varepsilon|Q(s)|$.
Moreover
$|1+\varepsilon-\varepsilon\kappa s|\le1+\varepsilon-\varepsilon\kappa s+2\varepsilon\kappa s\,[\varepsilon\kappa s>1]$.
Summing over $s\le D$ gives
\[
  \varphi\le(1+\varepsilon)\varphi-\varepsilon\kappa M
  -\varepsilon\Phi_{r,D}^{<c}(Q)+\varepsilon\Phi_{r,D}^{>c}(Q)
  +2\varepsilon\kappa\!\!\sum_{s>1/(\varepsilon\kappa),\ s\le D}\!\!s|P(s)|r^{-s}.
\]
Divide by $\varepsilon$ and let $\varepsilon\to0$.  The last sum tends to $0$
(it is empty for small $\varepsilon$ when $D<\infty$), so
\[
  \kappa M+\Phi_{r,D}^{<c}(Q)\le\varphi+\Phi_{r,D}^{>c}(Q).
\]

\emph{Step 4.}  By Step~2 the last inequality gives
$\kappa M_<+\kappa M_\ge(1-1/r)\le\varphi-\Phi_{r,D}^{<c}(Q)$.  So
$\frac\kappa rM_\ge\le(\varphi-\Phi_{r,D}^{<c}(Q))/(r-1)$.  Since $Q(c)=0$,
\[
  \Phi_{r,D}(Q)=\Phi_{r,D}^{<c}(Q)+\Phi_{r,D}^{>c}(Q)
  \le\Phi_{r,D}^{<c}(Q)+\frac{\varphi-\Phi_{r,D}^{<c}(Q)}{r-1}
  =\frac{\varphi+(r-2)\Phi_{r,D}^{<c}(Q)}{r-1}.
\]
For $r\ge2$ use $\Phi_{r,D}^{<c}(Q)\le\varphi$.  For $r<2$ drop the
nonpositive term.
\end{proof}

Only the optimality of $W$ in a single direction is used: the difference
$r_W-r_{\mathrm{ins}(W,c)}$ minus the multiple of $s\,r_W$ that cancels its
coefficient of degree $|W|+1$.

\begin{theorem}[Every row against the top row, with equality at roots at least $2$]\label{thm:allrowstwo}
Let $r>1$, $a=1/(r-1)$ and $n\ge1$.  For every finite $S\subset\N_{>0}$,
\[
  \mu_r\bigl(S,|S|+n\bigr)\le\frac{\max\{1,a\}^{|S|}}{\beta_r(n)} .
\]
Consequently $V_r(L,k)\ge\min\{1,r-1\}^{k-1}\beta_r(L-k+1)$ for all
$1\le k\le L$.  In particular, for every $r\ge2$, every $n\ge1$ and every
$L\ge n$,
\[
  V_r(L,L-n+1)=\beta_r(n).
\]
\end{theorem}

\begin{proof}
Put $\lambda=\max\{1,a\}$.  We show by induction on $|S|$ that
\[
  \mu_{r,D}\bigl(S,|S|+n\bigr)\le\lambda^{|S|}\mu_{r,D}(\varnothing,n)
  \qquad\text{for every integer }D\ge|S|+n+\max(S\cup\{0\}).
\]
For $S=\varnothing$ there is nothing to show.  Let $S\ne\varnothing$, and put
$\sigma=\max S$, $T=S\setminus\{\sigma\}$ and $t=\max(T\cup\{0\})$.  Then $D$
is admissible for $T$ as well.  By Lemma~\ref{lem:truncvertex}, some $W$ with
$T\subseteq W\subseteq[1,D]$ and $|W|=|T|+n-1$ has
$\Phi_{r,D}(r_W)=\mu_{r,D}(T,|T|+n)$.  Let $h=h_t(W)$ and
$W'=\mathrm{ins}(W,h)$.  The proof of Theorem~\ref{thm:crossing} uses only
the pointwise inequality $|q(s)|\le\lambda'|q_1(s)|+(1-\lambda')|q_2(s)|$ of
Lemma~\ref{lem:crossing}(a).  So it applies to the truncated sums, and gives
\[
  \mu_{r,D}\bigl(T\cup\{\sigma\},|W|+2\bigr)\le
  \max\bigl\{\Phi_{r,D}(r_W),\ \Phi_{r,D}(r_{W'})\bigr\}.
\]
Here $|W|+2=|S|+n$.  Lemma~\ref{lem:optins} applies with this $T$ and $W$
and with $c=h$, since $h\notin W$ and $h>t$.  It gives
$\Phi_{r,D}(r_{W'})\le\lambda\,\Phi_{r,D}(r_W)$.  Hence
$\mu_{r,D}(S,|S|+n)\le\lambda\,\mu_{r,D}(T,|T|+n)$, and the induction
hypothesis concludes.

Truncated sums are at most full ones, so
$\mu_{r,D}(\varnothing,n)\le\mu_r(\varnothing,n)=1/\beta_r(n)$.  The proof of
Proposition~\ref{prop:rowvalue}(b) shows that $\mu_{r,D}(S,L)\to\mu_r(S,L)$
as $D\to\infty$.  This gives the first claim.  Now let $1\le k\le L$ and
$n=L-k+1$.  Proposition~\ref{prop:rowvalue}(b) gives
$V_r(L,k)=\inf_{|S|=k-1}1/\mu_r(S,L)\ge\lambda^{-(k-1)}\beta_r(n)$, and
$\lambda^{-1}=\min\{1,r-1\}$.  For $r\ge2$ this is $\beta_r(n)$, and
Proposition~\ref{prop:rowvalue}(c) gives the reverse inequality.
\end{proof}

For $r\ge2$, Theorem~\ref{thm:prefixrows} gives
$V_r(L,k)\le\min_{i\le k}1/\nu_i(n)\le1/\nu_1(n)=\beta_r(n)$, so the theorem
also gives $\min_{i\le k}1/\nu_i(n)=\beta_r(n)$ and proves the prefix
identity~\eqref{eq:prefixid} there.  So the rows at a single
root $r\ge2$ are
\[
  V_r(L,k)=\beta_r(L-k+1)\qquad(1\le k\le L),
\]
the top row of a multiple of $(x-r)^{L-k+1}$.  At $r=2$ this sharpens
\cite[Theorem~4.13]{coefficient-mass} to $b_k(F)\ge\beta_2(L-k+1)$.

What remains open at a single root $r\ge2$ is the value of $\beta_r(n)$
itself.

For $1<r<2$ the theorem gives $V_r(L,k)\ge(r-1)^{k-1}\beta_r(L-k+1)$.  By
Corollary~\ref{cor:lastrowr} this is an equality for $k=L$, where both sides
are $(r-1)^L$.  With Theorem~\ref{thm:prefixrows} it shows that the
value $\min_{i\le k}1/\nu_i(n)$ of~\eqref{eq:prefixid} is at least
$(r-1)^{k-1}\beta_r(n)$.  Section~\ref{sec:onepoly} proves~\eqref{eq:prefixid}
for $1<r<2$ whenever $\beta_r(n)\le1$, and Section~\ref{sec:longdiag}
whenever $k$ is large compared with $n$.

For every $r>1$ and all $i,n\ge1$ the prefix values satisfy
\begin{equation}\label{eq:prefixup}
  \nu_{i+1}(n)\le a\,\nu_i(n),
\end{equation}
by Lemma~\ref{lem:optins}.  Take
$T=[1,i-1]$, a set $W$ given by Lemma~\ref{lem:truncvertex} for
$\mu_{r,D}([1,i-1],i-1+n)$, and $c=h_{i-1}(W)$.  Then $W\supseteq[1,c-1]$,
so $\Phi_{r,D}^{<c}(r_{\mathrm{ins}(W,c)})=0$, and
$\mathrm{ins}(W,c)\supseteq[1,i]$.  The lemma gives
$\mu_{r,D}([1,i],i+n)\le a\,\mu_{r,D}([1,i-1],i-1+n)$, and letting
$D\to\infty$ gives~\eqref{eq:prefixup}.

\section{Rows below the root \texorpdfstring{$2$}{2}}\label{sec:belowtwo}

For $1<r<2$ the rows are not all top rows, and Theorem~\ref{thm:allrowstwo}
gives only $V_r(L,k)\ge(r-1)^{k-1}\beta_r(n)$.  The prefix identity~\eqref{eq:prefixid}
would give $\min_{i\le k}1/\nu_i(n)$.  This section proves it in two regimes: for
every $k$ when $\beta_r(n)\le1$ (Section~\ref{sec:onepoly}), and for every
$n$ once $k$ is large compared with $n$ (Section~\ref{sec:longdiag}).  In
both the minimum is at $i=k$.  Section~\ref{sec:secondrow} treats the
second row for every $n$ and shows that the identity can fail, also
when $\nu_2(n)>\nu_1(n)$, and
Section~\ref{sec:finiterows} reduces every row to finitely many exempted
sets and shows that it can fail in the third row of a diagonal whose second
row satisfies it.

\subsection{One polynomial for all exempted sets}\label{sec:onepoly}

This subsection proves~\eqref{eq:prefixid} for $1<r<2$ whenever $\beta_r(n)\le1$,
for every $k$ (Theorem~\ref{thm:onepolyrows}).  There the minimum is at
$i=k$: the prefix $[1,k-1]$ is a worst exempted set.

It rests on two facts: one polynomial $p$ serves all exempted sets at
once, at the cost of the largest of its prefix sums
(Theorem~\ref{thm:onepoly}); and these prefix sums grow at least by the
factor $\nu_1(n)=1/\beta_r(n)$ from one index to the next
(Lemma~\ref{lem:prefixshift}).  For the first,
Lemma~\ref{lem:crossing}(b) pushes every exempted position down into an
initial block while the fixed factor $p$ rides along.  Nothing is inserted,
and the only optimality used is that of $p$ for the last index.

Put $\omega_i(s)=\binom{s-1}{i-1}r^{-s}$ for $i\ge1$ (zero for $s<i$),
let $\mathcal Q_n$ be the set of real polynomials $p$ with
$\deg p\le n-1$ and $p(0)=1$, and put
\[
  A_i(p)=\sum_{s\ge1}\omega_i(s)\,|p(s)|=\Phi_r\bigl(r_{[1,i-1]}\,p\bigr)
  \qquad(i\ge1),
\]
for real polynomials $p$, so that $\nu_i(n)=\inf_{p\in\mathcal Q_n}A_i(p)$.
The second expression holds because $|r_{[1,i-1]}(s)|=\binom{s-1}{i-1}$
for $s\ge1$.  For finite $S\subset\N_{>0}$ let $\ell(S)$ be its leading
run (Section~\ref{sec:rowvalues}), the largest $\ell\ge0$ with
$[1,\ell]\subseteq S$.

\begin{lemma}[Pushing exempted positions into the prefix]\label{lem:prefixpush}
Let $r>1$, let $p$ be a real polynomial and let $S\subset\N_{>0}$ be finite.
Then
\[
  \Phi_r(r_S\,p)\le\max_{1\le i\le|S|+1}A_i(p).
\]
\end{lemma}

\begin{proof}
Induct on $|S|$ and, for fixed $|S|$, on $d(S)=|S|-\ell(S)$.  If $d(S)=0$,
then $S=[1,|S|]$ and $\Phi_r(r_Sp)=A_{|S|+1}(p)$.  Let $d(S)>0$, put
$\ell=\ell(S)$ and $\sigma=\max S$.  Then $\ell+1\notin S$, so
$\sigma\ge\ell+2$; otherwise $S\subseteq[1,\ell]$ and $d(S)=0$.  Put
$U=S\setminus\{\sigma\}$ and $Q=r_Up$, so that $r_Sp=Q\cdot(1-s/\sigma)$.
Lemma~\ref{lem:crossing}(b) with $y=\ell+1\le\sigma$ gives
\[
  \Phi_r(r_Sp)\le\max\bigl\{\Phi_r(r_Up),\ \Phi_r(r_{U\cup\{\ell+1\}}\,p)\bigr\},
\]
because $\ell+1\notin U$.  Here $|U|=|S|-1$, so the induction hypothesis
bounds the first term by $\max_{i\le|S|}A_i(p)$.  The set
$S'=U\cup\{\ell+1\}$ has $|S'|=|S|$ and $[1,\ell+1]\subseteq S'$, so
$d(S')<d(S)$, and the induction hypothesis bounds the second term by
$\max_{i\le|S|+1}A_i(p)$.
\end{proof}

\begin{theorem}[One polynomial for all exempted sets]\label{thm:onepoly}
Let $r>1$, $1\le k\le L$, $n=L-k+1$ and $p\in\mathcal Q_n$.  Every monic
multiple $F$ of $(x-r)^L$ has
\[
  b_k(F)\ \ge\ \frac1{\max_{1\le i\le k}A_i(p)} .
\]
Hence, with $\Lambda_k(n)=\inf_{p\in\mathcal Q_n}\max_{1\le i\le k}A_i(p)$,
\[
  \frac1{\Omega_k(n)}\ \le\ \frac1{\Lambda_k(n)}\ \le\ V_r(L,k)\ \le\
  \min_{1\le i\le k}\frac1{\nu_i(n)} .
\]
\end{theorem}

\begin{proof}
Let $S$ be the exempted set of Proposition~\ref{prop:rowvalue}(a), so
$|S|=k-1$.  The polynomial $q=r_Sp$ has $\deg q\le k-1+n-1=L-1$, $q(0)=1$
and $q|_S=0$.  So $b_k(F)\,\Phi_r(q)\ge1$, and Lemma~\ref{lem:prefixpush}
bounds $\Phi_r(q)$.  The upper bound is Theorem~\ref{thm:prefixrows}.  For
the first inequality, $W_k(s)=\max_{1\le i\le k}\binom{s-1}{i-1}$ in the
notation of Theorem~\ref{thm:prefixrows}, so
$\sum_sW_k(s)r^{-s}|p(s)|\ge\max_{i\le k}A_i(p)$ for every $p$.
\end{proof}

The theorem moves the maximum over $i$ out of the sum in
Theorem~\ref{thm:prefixrows}, and this loses nothing when one $p$ is nearly
optimal for all prefix problems at once.  The prefix identity~\eqref{eq:prefixid} holds at $(r,L,k)$ whenever
$\Lambda_k(n)=\max_{i\le k}\nu_i(n)$.  For $n=1$ we have
$\mathcal Q_1=\{1\}$ and $A_i(1)=a^i$.  So the theorem gives
$V_r(L,L)\ge r-1$ for $r\ge2$ and $V_r(L,L)\ge(r-1)^L$ for $r\le2$, the
lower bounds of Corollary~\ref{cor:lastrowr}; its proof is this argument,
through Lemma~\ref{lem:prefixpush} with $p=1$.

\begin{lemma}[Shifting the prefix index]\label{lem:prefixshift}
Let $r>1$, $n\ge1$, $i,m\ge1$, and let $p$ be real with $\deg p\le n-1$.
Then
\[
  A_{i+m}(p)\ \ge\ \nu_m(n)\,A_i(p).
\]
Consequently $\nu_{i+m}(n)\ge\nu_i(n)\,\nu_m(n)$, and
$\nu_i(n)\ge\nu_1(n)^i=\beta_r(n)^{-i}$.
\end{lemma}

\begin{proof}
Counting the $(i+m-1)$-element subsets of $[1,s-1]$ by their $i$-th
smallest element $u$ gives
$\sum_{u=1}^{s-1}\binom{u-1}{i-1}\binom{s-u-1}{m-1}=\binom{s-1}{i+m-1}$, that is,
$\omega_{i+m}(s)=\sum_{u+t=s,\ u,t\ge1}\omega_i(u)\,\omega_m(t)$.  All terms
being nonnegative,
\[
  A_{i+m}(p)=\sum_{u\ge1}\omega_i(u)\sum_{t\ge1}\omega_m(t)\,|p(u+t)| .
\]
If $p(u)\ne0$, the polynomial $t\mapsto p(u+t)/p(u)$ lies in $\mathcal Q_n$,
so the inner sum is at least $\nu_m(n)|p(u)|$; if $p(u)=0$ this is trivial.
This proves the inequality.  Taking the infimum over $p\in\mathcal Q_n$
gives $\nu_{i+m}(n)\ge\nu_m(n)\nu_i(n)$, and $\nu_1(n)=1/\beta_r(n)$.
\end{proof}

Together with~\eqref{eq:prefixup},
$\nu_1(n)\,\nu_i(n)\le\nu_{i+1}(n)\le a\,\nu_i(n)$, and both inequalities are
equalities at $n=1$, where $\nu_i(1)=a^i$.

\begin{theorem}[All rows when the top row is at most one]\label{thm:onepolyrows}
Let $r>1$ and $n\ge1$ with $\beta_r(n)\le1$.  Then
$1\le\nu_1(n)\le\nu_2(n)\le\cdots$, and for every $k\ge1$ and $L=n+k-1$,
\[
  V_r(L,k)=\frac1{\nu_k(n)}=\min_{1\le i\le k}\frac1{\nu_i(n)} .
\]
So the prefix identity~\eqref{eq:prefixid} holds at $(r,n+k-1,k)$
for every $k$, with the minimum at $i=k$.
\end{theorem}

\begin{proof}
By Lemma~\ref{lem:prefixshift} with $m=1$ and $\nu_1(n)=1/\beta_r(n)\ge1$,
$A_{i+1}(p)\ge A_i(p)$ for every $p\in\mathcal Q_n$ and every $i$, and
$\nu_{i+1}(n)\ge\nu_i(n)$.  Given $\varepsilon>0$ choose $p\in\mathcal Q_n$
with $A_k(p)\le\nu_k(n)+\varepsilon$.  Then
$\max_{i\le k}A_i(p)=A_k(p)$, and Theorem~\ref{thm:onepoly} gives
$V_r(L,k)\ge1/(\nu_k(n)+\varepsilon)$.  The upper bound
$V_r(L,k)\le\min_i1/\nu_i(n)\le1/\nu_k(n)$ is Theorem~\ref{thm:prefixrows}.
\end{proof}

The hypothesis $\beta_r(n)\le1$ holds for $r$ in an interval with left end
$1$, which is $(1,2]$ for $n=1$, shrinks as $n$ grows, and is nonempty for
every $n$ (Lemma~\ref{lem:topnearone}).  Indeed, since
$\nu_1(n+1)\le\nu_1(n)$, we have $\beta_r(n)\ge\beta_r(1)=r-1$.  So the
hypothesis forces $r\le2$, and at $r=2$ it holds only for $n=1$, since
$\beta_2(n)\ge n$ by \cite[Theorem~4.13]{coefficient-mass}.  For each
$q\in\mathcal Q_n$ the sum $\Phi_r(q)$ decreases in $r$, term by term, so
$\nu_1(n)$ does not increase with $r$, and the set
$\{r>1:\beta_r(n)\le1\}$ is an interval with left end $1$.

\begin{lemma}[The top row near the root $1$]\label{lem:topnearone}
Let $n\ge1$, $1<r\le2$ and $C_n=\prod_{i=1}^n(2i+1)$.  Then
$\beta_r(n)\le e^{4n+2}C_n\,(r-1)$.  In particular Theorem~\ref{thm:onepolyrows}
applies whenever $r-1\le e^{-4n-2}/C_n$.
\end{lemma}

\begin{proof}
Let $q\in\mathcal Q_n$ and $m=\lceil1/(r-1)\rceil$, so $1\le(r-1)m\le2$.
For $1\le j\le n$ let $B_j$ be the $m$ integers in $[2jm,2jm+m)$.  Fix
$s_j\in B_j$.  Lagrange interpolation at $s_1,\ldots,s_n$ gives
$1=q(0)=\sum_j\lambda_jq(s_j)$ with
$\lambda_j=\prod_{i\ne j}s_i/(s_i-s_j)$.  Here $s_i<(2i+1)m$ and
$|s_i-s_j|>(2|i-j|-1)m\ge m$ for $i\ne j$, so $|\lambda_j|\le C_n$ and
$1\le C_n\sum_j|q(s_j)|$.  Averaging over the $m^n$ choices of
$(s_1,\ldots,s_n)$ gives $m\le C_n\sum_{s\in B}|q(s)|$, where
$B=B_1\cup\cdots\cup B_n\subseteq[1,(2n+1)m]$.  For $s\le(2n+1)m$,
$r^{-s}\ge e^{-(r-1)(2n+1)m}\ge e^{-4n-2}$, since $\log r\le r-1$.  Hence
$\Phi_r(q)\ge e^{-4n-2}m/C_n\ge e^{-4n-2}/(C_n(r-1))$, and taking the
infimum over $q$ gives $\nu_1(n)\ge e^{-4n-2}/(C_n(r-1))$, which is the
claim since $\beta_r(n)=1/\nu_1(n)$.
\end{proof}

Beyond $\beta_r(n)\le1$, where the argument of Theorem~\ref{thm:onepolyrows}
no longer shows that one polynomial serves all prefixes,
we do not know in general when the prefix identity holds;
Section~\ref{sec:secondrow} reduces the case $k=2$ to finitely many exempted
positions, and Section~\ref{sec:finiterows} every $k$ to finitely many
exempted sets.  Section~\ref{sec:longdiag} shows that a
single polynomial does serve all prefixes once $k$ is large compared with
$n$.

\subsection{Long diagonals}\label{sec:longdiag}

Theorem~\ref{thm:onepoly} gives the value $1/\nu_k(n)$ as soon
as one polynomial that is nearly optimal for $\nu_k(n)$ has its $k$-th
prefix sum $A_k$ at least as large as the earlier ones.
Theorem~\ref{thm:onepolyrows} obtains this from $\beta_r(n)\le1$.  This
subsection obtains it for every $1<r<2$, with no condition on $\beta_r(n)$,
once $k$ is large compared with $n$ (Theorem~\ref{thm:longdiag} and
Proposition~\ref{prop:longdiagexplicit}).  So on every diagonal
$L-k=n-1$ the prefix identity can fail in at most finitely many rows
(Corollary~\ref{cor:longdiagfinite}).

The earlier weights $\omega_i$, $i<k$, can exceed $\omega_k$ only on
$[1,2k-3]$, where a suitable polynomial for $\nu_k(n)$ is bounded by $1$;
when this excess is small, the $k$-th prefix sum of that polynomial is the
largest.  By
Lemma~\ref{lem:intzeros} the polynomials for
$\nu_k(n)$ may be taken with all zeros at least $k$, and such a polynomial is
bounded by $1$ on $[1,2k-3]$.  There the weights $\omega_i$, $i<k$, exceed
$\omega_k$ by an amount that is exponentially small compared with $\nu_k(n)$
when $r<2$ and $k$ is large compared with $n$.  From $2k-2$ on, $\omega_k$ dominates every $\omega_i$, $i<k$,
with a margin that is a fixed fraction of $\omega_k$ near the mode of
$\omega_k$, where every polynomial in $\mathcal Q_n$ has some mass.

For $k\ge2$ put
\[
  w_k(s)=\omega_k(s)-\omega_{k-1}(s),\qquad
  \tau_k(n)=\inf_{p\in\mathcal Q_n}\sum_{s\ge2k-2}w_k(s)\,|p(s)|,
\]
\[
  N_k(r)=\max_{1\le i<k}\sum_{s=1}^{2k-3}\bigl(\omega_i(s)-\omega_k(s)\bigr)^+,
\]
where $x^+=\max\{x,0\}$ and $\omega_i(s)=0$ for $s<i$.  For $s\ge k$ we have
$\omega_{k-1}(s)/\omega_k(s)=(k-1)/(s-k+1)$, so
\begin{equation}\label{eq:wkratio}
  w_k(s)=\frac{s-2k+2}{s-k+1}\,\omega_k(s)\qquad(s\ge k).
\end{equation}
In particular $w_k\ge0$ on $[2k-2,\infty)$, and $\tau_k(n)\ge0$.

\begin{lemma}[Far zeros keep the last prefix sum largest]\label{lem:farprefix}
Let $r>1$, $k\ge2$ and $n\ge1$.  Let $Y$ be a set of $n-1$ integers, each at
least $k$, and $p=r_Y$.  Then for $1\le i<k$,
\[
  A_k(p)-A_i(p)\ \ge\ \tau_k(n)-N_k(r).
\]
\end{lemma}

\begin{proof}
Let $1\le s\le2k-3$.  Every $y\in Y$ has $2y\ge2k>s$, so $|1-s/y|\le1$ and
$|p(s)|\le1$.  Hence
$(\omega_k(s)-\omega_i(s))|p(s)|\ge-(\omega_i(s)-\omega_k(s))^+$, and the
sum of these terms over $s\le2k-3$ is at least $-N_k(r)$.

Let $s\ge2k-2$.  The coefficients $\binom{s-1}j$ increase in $j$ for
$j\le(s-1)/2$, and $i-1\le k-2\le(s-1)/2$, so $\omega_i(s)\le\omega_{k-1}(s)$
and $(\omega_k(s)-\omega_i(s))|p(s)|\ge w_k(s)|p(s)|$.  Since $p\in\mathcal Q_n$,
the sum of these terms over $s\ge2k-2$ is at least $\tau_k(n)$.
\end{proof}

\begin{theorem}[Long diagonals]\label{thm:longdiag}
Let $r>1$, $n\ge1$, $k\ge2$ and $L=n+k-1$.  If $N_k(r)\le\tau_k(n)$, then
$\nu_i(n)\le\nu_k(n)$ for $1\le i\le k$, and
\[
  V_r(L,k)=\frac1{\nu_k(n)}=\min_{1\le i\le k}\frac1{\nu_i(n)} .
\]
So the prefix identity~\eqref{eq:prefixid} holds at
$(r,L,k)$, with the minimum at $i=k$.
\end{theorem}

\begin{proof}
Let $\varepsilon>0$.  Lemma~\ref{lem:intzeros} with $\omega=\omega_k$ gives a
set $Y$ of $n-1$ integers, each at least $k$, such that $p=r_Y$ has
$A_k(p)\le\nu_k(n)+\varepsilon$.  By Lemma~\ref{lem:farprefix} and the
hypothesis, $A_i(p)\le A_k(p)$ for $i<k$.  Hence
$\nu_i(n)\le A_i(p)\le\nu_k(n)+\varepsilon$ for $i\le k$, and
$\max_{i\le k}A_i(p)=A_k(p)$.  Theorem~\ref{thm:onepoly} gives
$V_r(L,k)\ge1/(\nu_k(n)+\varepsilon)$.  Let $\varepsilon\to0$.
Theorem~\ref{thm:prefixrows} gives
$V_r(L,k)\le\min_i1/\nu_i(n)\le1/\nu_k(n)$.
\end{proof}

The hypothesis compares a quantity that is exponentially small relative to
$a^k$ for $r<2$ with one that is not.  The next proposition makes this explicit.  It uses
the mode of $\omega_k$ and Lagrange interpolation, as in
Lemma~\ref{lem:nbmass}.

\begin{proposition}[An explicit form]\label{prop:longdiagexplicit}
Let $1<r<2$, $\theta=4(r-1)/r^2$ and $n\ge1$.  For $k\ge2$ put
$M_k=\lfloor(k-1)r/(r-1)\rfloor+1$ and
\[
  \Xi(k)=\theta^k\,(2r)^{n-1}\binom{M_k+n-1}{n-1}
  \Bigl(\frac{4\sqrt{kr}}{r-1}+1\Bigr).
\]
\begin{enumerate}
\item[(a)] If $\Xi(k)\le3(2-r)^2/r^2$, then $N_k(r)\le\tau_k(n)$, so
Theorem~\ref{thm:longdiag} applies.
\item[(b)] Put $d=\lceil r/(r-1)\rceil$ and
$\gamma_k=\theta\,(1+(n-1)/(M_k+1))^d\sqrt{1+1/k}$.  Then $\Xi(k+1)\le\gamma_k\Xi(k)$,
and $\gamma_k$ does not increase with $k$.  So if $\Xi(k_1)\le3(2-r)^2/r^2$ and
$\gamma_{k_1}\le1$, then (a) applies to every $k\ge k_1$.
\item[(c)] Let $\alpha\ge1$ satisfy
$\rho(\alpha):=\theta\,\bigl(2re(\alpha r/(r-1)+1)\bigr)^{1/\alpha}<1$.  If
$n-1\le k/\alpha$, then
$\Xi(k)\le\rho(\alpha)^k(4\sqrt{kr}/(r-1)+1)$.  Consequently there is $K$ such
that $V_r(n+k-1,k)=1/\nu_k(n)$ for every $n\ge1$ and every
$k\ge\max\{K,\alpha(n-1)\}$.
\end{enumerate}
Since $\bigl(2re(\alpha r/(r-1)+1)\bigr)^{1/\alpha}\to1$ as $\alpha\to\infty$
and $\theta<1$, admissible $\alpha$ exist for every $1<r<2$.
\end{proposition}

\begin{proof}
(a) First, $(\omega_i(s)-\omega_k(s))^+\le\omega_i(s)\le
\binom{s-1}{\lfloor(s-1)/2\rfloor}r^{-s}$, so $N_k(r)\le E_k(r)$, and the
proof of Corollary~\ref{cor:belowtwo} gives
$E_k(r)\le\frac{r^2}{4(2-r)}(4/r^2)^k=\frac{r^2}{4(2-r)}\,\theta^ka^k$.

Next let $p\in\mathcal Q_n$.  The ratio $\omega_k(s+1)/\omega_k(s)=s/(r(s-k+1))$
is at least $1/r$, and it is at least $1$ exactly when $s\le(k-1)r/(r-1)$.  So
$M_k$ maximizes $\omega_k$: it is a mode of the law $P$ of
Lemma~\ref{lem:nbmass}, where $\omega_k=a^kP$.  That lemma gives
$\omega_k(M_k)\ge a^k\cdot3/(4(4\sqrt{kr}/(r-1)+1))$, and
$\omega_k(M_k+l)\ge r^{-l}\omega_k(M_k)$ for $l\ge0$.  Since
$M_k>(k-1)r/(r-1)\ge2k-2$, the nodes $t_l=M_k+l$, $0\le l<n$, lie in
$[2k-2,\infty)$.  By~\eqref{eq:wkratio} the ratio $w_k(s)/\omega_k(s)=1-(k-1)/(s-k+1)$
increases in $s$, and at $s=t_0>(k-1)r/(r-1)$ it exceeds $1-(r-1)=2-r$.  So
$w_k(t_l)\ge(2-r)\,r^{-(n-1)}\omega_k(M_k)$ for every $l$.  Finally
$1=p(0)=\sum_l\ell_lp(t_l)$ with the Lagrange coefficients
$|\ell_l|=\prod_{l'\ne l}t_{l'}/|l-l'|\le\prod_{j=1}^{n-1}(M_k+j)/(l!\,(n-1-l)!)$,
so $\sum_l|\ell_l|\le2^{n-1}\binom{M_k+n-1}{n-1}$ and
$\max_l|p(t_l)|\ge1/(2^{n-1}\binom{M_k+n-1}{n-1})$.  Altogether
\[
  \sum_{s\ge2k-2}w_k(s)|p(s)|\ge\sum_lw_k(t_l)|p(t_l)|
  \ge\frac{3(2-r)\,a^k}{4\,(4\sqrt{kr}/(r-1)+1)\,(2r)^{n-1}\binom{M_k+n-1}{n-1}} .
\]
This lower bound for $\tau_k(n)$ is at least the upper bound for $E_k(r)$
exactly when $\Xi(k)\le3(2-r)^2/r^2$.

(b) We have $M_{k+1}-M_k\le d$, so
$\binom{M_{k+1}+n-1}{n-1}/\binom{M_k+n-1}{n-1}=\prod_{j=M_k+1}^{M_{k+1}}(j+n-1)/j
\le(1+(n-1)/(M_k+1))^d$.  Also
$(4\sqrt{(k+1)r}/(r-1)+1)/(4\sqrt{kr}/(r-1)+1)\le\sqrt{(k+1)/k}$.  This gives
$\Xi(k+1)\le\gamma_k\Xi(k)$.  Since $M_k$ does not decrease, neither does
$1/\gamma_k$, and induction on $k$ gives the rest.

(c) For $n=1$ the claim is $\theta\le\rho(\alpha)$.  Let $m=n-1\ge1$.  Then
$\binom{M_k+m}m\le(e(M_k+m)/m)^m$, and the right side increases with $m$,
since the derivative of $m\log(e(M_k+m)/m)$ is $\log(1+M_k/m)+m/(M_k+m)>0$.
With $m\le k/\alpha$ and $M_k\le(k-1)r/(r-1)+1\le kr/(r-1)$ this gives
$\binom{M_k+m}m\le(e(\alpha r/(r-1)+1))^{k/\alpha}$.  Also
$(2r)^m\le(2r)^{k/\alpha}$, which proves the bound on $\Xi(k)$.  Its right side
tends to $0$ as $k\to\infty$ and decreases once $k\ge1/(2\log(1/\rho(\alpha)))$.
So some $K$ has $\rho(\alpha)^k(4\sqrt{kr}/(r-1)+1)\le3(2-r)^2/r^2$ for all
$k\ge K$, and (a) and Theorem~\ref{thm:longdiag} apply whenever also
$k\ge\alpha(n-1)$.
\end{proof}

\begin{corollary}[Every diagonal is eventually determined]\label{cor:longdiagfinite}
Let $1<r<2$ and $n\ge1$.  Then $V_r(L,L-n+1)=1/\nu_{L-n+1}(n)$ for all but
finitely many $L\ge n$.  Moreover the sequence $\nu_k(n)$ is eventually
nondecreasing in $k$, and $\nu_k(n)=\max_{i\le k}\nu_i(n)$ for all large $k$.
\end{corollary}

\begin{proof}
Apply Proposition~\ref{prop:longdiagexplicit}(c) with $k=L-n+1\ge\max\{K,\alpha(n-1)\}$.
Theorem~\ref{thm:longdiag} also gives $\nu_i(n)\le\nu_k(n)$ for all
$i\le k$ and all such $k$.
\end{proof}

This sharpens Corollary~\ref{cor:belowtwo}, which gives
$V_r(L,k)=(1+o(1))/\nu_k(n)$ as $k\to\infty$ with $n$ fixed: the factor
$1+o(1)$ is exactly $1$ once $k\ge\max\{K,\alpha(n-1)\}$, where $K$ does not
depend on $n$.
What remains open for $1<r<2$ are the rows with $\beta_r(n)>1$ and $k$
below these thresholds.  The next subsection treats the first of them,
$k=2$, and Section~\ref{sec:finiterows} reduces each of them to finitely
many exempted sets.

\subsection{The second row}\label{sec:secondrow}

This subsection treats the second row, $k=2$, with no condition on
$\beta_r(n)$, through the values $\eta_\sigma(n)=\mu_r(\{\sigma\},n+1)$ of
single exempted positions $\sigma$.  A polynomial vanishing at $\sigma$ is
$(1-s/\sigma)p$ with $p\in\mathcal Q_n$, so
\[
  \eta_\sigma(n)=\mu_r(\{\sigma\},n+1)
  =\inf_{p\in\mathcal Q_n}\sum_{s\ge1}\Bigl|1-\frac s\sigma\Bigr||p(s)|\,r^{-s}
  \qquad(\sigma\in\N_{>0}).
\]
Then $\eta_1(n)=\nu_2(n)$, and Proposition~\ref{prop:rowvalue}(b) gives
$V_r(n+1,2)=\inf_\sigma1/\eta_\sigma(n)$.  Far positions cost no more than
the limit $\nu_1(n)$, so the second row is the minimum of $1/\nu_1(n)$ and
finitely many position terms (Theorem~\ref{thm:secondrow}).  Each term is
$1/\eta_\sigma(n)=1/\Phi_r(Z)$ for a set $Z\ni\sigma$ of $n$ integers, and
such a $Z$ can be certified exactly
(Lemma~\ref{lem:vertexopt}).  At $r=\frac54$ and $r=\frac43$
this determines rows on both sides of the prefix identity
(Proposition~\ref{prop:secondrowexact}), and at $r=\frac{177}{167}$ it
shows that the identity can fail although $\nu_2(n)>\nu_1(n)$
(Proposition~\ref{prop:secondrowhyp}).

\begin{lemma}[Dropping a far exempted position]\label{lem:rowmono}
Let $r>1$ and $2\le k\le L$.  Then $V_r(L,k)\le V_r(L-1,k-1)$.  So
$V_r(n+k-1,k)$ does not increase with $k$, for each $n\ge1$.
\end{lemma}

\begin{proof}
Let $S'\subset\N_{>0}$ with $|S'|=k-2$, let $N>\max(S'\cup\{0\})$ and
$S_N=S'\cup\{N\}$.  A real $q$ admissible for $\mu_r(S_N,L)$ factors as
$q=(1-s/N)q_0$ with $q_0$ admissible for $\mu_r(S',L-1)$.  For
$1\le s\le M<N$ we have $1-s/N\ge1-M/N$, so
$\Phi_r(q)\ge(1-M/N)\,\mu_{r,M}(S',L-1)$.  Hence
$\liminf_N\mu_r(S_N,L)\ge\mu_{r,M}(S',L-1)$ for every $M$, and by the
convergence in the proof of Proposition~\ref{prop:rowvalue}(b) it is at
least $\mu_r(S',L-1)$.  Proposition~\ref{prop:rowvalue}(b) gives
$V_r(L,k)\le\inf_N1/\mu_r(S_N,L)\le1/\mu_r(S',L-1)$, and the infimum over
$S'$ is $V_r(L-1,k-1)$, again by Proposition~\ref{prop:rowvalue}(b).
\end{proof}

With $V_r(n+i-1,i)\le1/\mu_r([1,i-1],n+i-1)=1/\nu_i(n)$ this recovers the
upper bound of Theorem~\ref{thm:prefixrows}.  It also propagates failures
along a diagonal: if $V_r(n+1,2)<\min_{i\le k}1/\nu_i(n)$, then the prefix
identity fails at $(r,n+k-1,k)$.

For a real polynomial $p$ and real $\sigma>0$ put
\[
  \Psi_r(p)=\sum_{s\ge1}s\,|p(s)|\,r^{-s},\qquad
  R_p(\sigma)=\sum_{s>\sigma}(s-\sigma)\,|p(s)|\,r^{-s}.
\]
Since $|1-s/\sigma|=1-s/\sigma+2(s/\sigma-1)^+$, multiplying by
$|p(s)|r^{-s}$ and summing (all three series converge absolutely) gives
\[
  \Phi_r\bigl((1-s/\sigma)p\bigr)=\Phi_r(p)-\frac{\Psi_r(p)-2R_p(\sigma)}\sigma .
\]
If $p\ne0$, then $\Psi_r(p)\ge\Phi_r(p)>0$ and
$R_p(\sigma)\le\sum_{s>\sigma}s|p(s)|r^{-s}\to0$, so a zero far enough out
lowers $\Phi_r$ strictly.

\begin{lemma}[Certifying a vertex]\label{lem:vertexopt}
Let $r>1$, let $S\subseteq Z$ be finite subsets of $\N_{>0}$ and
$L=|Z|+1$.  For $y\in Z\setminus S$ put $e_y(s)=s\,r_{Z\setminus\{y\}}(s)$
and
\[
  G_y=\sum_{s\ge1,\ s\notin Z}\sgn\bigl(r_Z(s)\bigr)\,e_y(s)\,r^{-s}.
\]
Then $\Phi_r(Z)=\mu_r(S,L)$ if and only if $|G_y|\le|e_y(y)|\,r^{-y}$ for
every $y\in Z\setminus S$.  If $|G_y|<|e_y(y)|\,r^{-y}$ for every
$y\in Z\setminus S$, then $r_Z$ is the unique minimizer: every other real
$q$ with $\deg q\le L-1$, $q(0)=1$ and $q|_S=0$ has
$\Phi_r(q)>\Phi_r(Z)$.  Moreover, for every finite
$S\subset\N_{>0}$ and $L>|S|$ some $Z\supseteq S$ with $|Z|=L-1$ has
$\Phi_r(Z)=\mu_r(S,L)$.
\end{lemma}

\begin{proof}
The polynomials admissible for $\mu_r(S,L)$ form the affine space
$r_Z+E$, where $E$ is the space of real $d$ with $\deg d\le L-1$, $d(0)=0$
and $d|_S=0$, of dimension $L-1-|S|=|Z\setminus S|$.  Each $e_y$ lies in
$E$, and $e_y(y')=0\ne e_y(y)$ for $y\ne y'$ in $Z\setminus S$, so the $e_y$
form a basis and $d=\sum_yc_ye_y$ has $d(y)=c_ye_y(y)$.  The function
$\Phi_r$ is convex and finite on $r_Z+E$, so $r_Z$ is a minimizer if and
only if every one-sided derivative $\Phi_r'(r_Z;d)$ is nonnegative.  For
$s\notin Z$ we have $r_Z(s)\ne0$ and
$(|r_Z(s)+\varepsilon d(s)|-|r_Z(s)|)/\varepsilon\to\sgn(r_Z(s))\,d(s)$ as
$\varepsilon\downarrow0$; for $s\in S$ both terms vanish; for
$s=y\in Z\setminus S$ the quotient is $|d(y)|$.  The quotients are bounded
by $|d(s)|$ and $\sum_s|d(s)|r^{-s}<\infty$, so by dominated convergence
\[
  \Phi_r'(r_Z;d)=\sum_{y\in Z\setminus S}\bigl(c_yG_y+|c_y|\,|e_y(y)|\,r^{-y}\bigr).
\]
This is nonnegative for all $c$ if and only if it is for each $c$
supported on one $y$, with either sign, that is, if and only if
$|G_y|\le|e_y(y)|r^{-y}$ for every $y$.

If every inequality is strict and $d=\sum_yc_ye_y\ne0$, then some
$c_y\ne0$, and each term of $\Phi_r'(r_Z;d)$ is at least
$|c_y|(|e_y(y)|r^{-y}-|G_y|)\ge0$, with strict inequality when $c_y\ne0$;
so $\Phi_r'(r_Z;d)>0$.  By convexity the quotient
$(\Phi_r(r_Z+\varepsilon d)-\Phi_r(r_Z))/\varepsilon$ does not decrease in
$\varepsilon\in(0,1]$, so
$\Phi_r(r_Z+d)-\Phi_r(r_Z)\ge\Phi_r'(r_Z;d)>0$.  Every admissible
$q\ne r_Z$ is $r_Z+d$ with such a $d$.

For the last claim take integers $D_j\to\infty$ with
$D_j\ge L+\max(S\cup\{0\})$ and $W_j$ as in Lemma~\ref{lem:truncvertex} for
$\mu_{r,D_j}(S,L)$.  Passing to a subsequence, for each $i$ the $i$-th
smallest element of $W_j$ is eventually constant or tends to $\infty$.  A
factor $1-s/w$ with $w\to\infty$ tends to $1$, so $r_{W_j}\to r_Z$
coefficientwise, where $Z\supseteq S$, $|Z|\le L-1$, consists of the
eventually constant elements.  For each $M$,
$\sum_{s\le M}|r_Z(s)|r^{-s}=\lim_j\sum_{s\le M}|r_{W_j}(s)|r^{-s}$, and for
$D_j\ge M$ each partial sum $\sum_{s\le M}|r_{W_j}(s)|r^{-s}$ is at most
$\Phi_{r,D_j}(r_{W_j})=\mu_{r,D_j}(S,L)\le\mu_r(S,L)$.  So
$\Phi_r(Z)\le\mu_r(S,L)$, with equality because $r_Z$ is admissible.  If
$|Z|<L-1$, then $(1-s/\sigma)r_Z$ is admissible for $\mu_r(S,L)$ for every
$\sigma>0$, and for large $\sigma$ the identity before the lemma gives
$\Phi_r((1-s/\sigma)r_Z)<\Phi_r(Z)=\mu_r(S,L)$, a contradiction.  So
$|Z|=L-1$.
\end{proof}

Beyond $\max Z$ the sign of $r_Z$ is $(-1)^{|Z|}$, and
$\sum_{s\ge0}P(s)x^s=\sum_j(\Delta^jP)(0)\,x^j/(1-x)^{j+1}$ for a
polynomial $P$, where $(\Delta P)(s)=P(s+1)-P(s)$ is the forward
difference.  So for rational $r$ both $\Phi_r(Z)$ and the $G_y$ are
rational numbers computable exactly, and for a given $Z$ the test of the
lemma is finite.  By the last claim, every $\mu_r(S,L)$, in particular every
$\eta_\sigma(n)$, is attained at some such $Z$; the lemma does not say how
to find it.


\begin{theorem}[The second row is a finite minimum]\label{thm:secondrow}
Let $r>1$ and $n\ge1$, and put $\eta_\sigma(n)=\mu_r(\{\sigma\},n+1)$ for
$\sigma\in\N_{>0}$.
\begin{enumerate}
\item[(a)] For every $p\in\mathcal Q_n$ and $\sigma\in\N_{>0}$,
\[
  \eta_\sigma(n)\le\Phi_r(p)-\frac{\Psi_r(p)-2R_p(\sigma)}\sigma .
\]
\item[(b)] Some $p_*\in\mathcal Q_n$ has $\Phi_r(p_*)=\nu_1(n)$.  Fix any
such $p_*$, and let $\sigma_*$ be the least $\sigma\in\N_{>0}$ with
$2R_{p_*}(\sigma)\le\Psi_r(p_*)$.  Then $\eta_\sigma(n)\le\nu_1(n)$ for
$\sigma\ge\sigma_*$, with strict
inequality for $\sigma>\sigma_*$, and
\[
  V_r(n+1,2)=\min\Bigl\{\frac1{\nu_1(n)},\
  \min_{1\le\sigma<\sigma_*}\frac1{\eta_\sigma(n)}\Bigr\},
\]
where the inner minimum is omitted if $\sigma_*=1$.
\item[(c)] The prefix identity~\eqref{eq:prefixid} fails at $(r,n+1,2)$
if and only if $\eta_\sigma(n)>\max\{\nu_1(n),\nu_2(n)\}$ for some
$2\le\sigma<\sigma_*$.
\end{enumerate}
Parts (b) and (c) hold for every choice of $p_*$.
\end{theorem}

\begin{proof}
(a) The polynomial $(1-s/\sigma)p$ is admissible for $\eta_\sigma(n)$, and
the identity before Lemma~\ref{lem:vertexopt} gives its value $\Phi_r$.

(b) For $n=1$, $\mathcal Q_1=\{1\}$.  For $n\ge2$, the values
$p(1),\ldots,p(n-1)$ are affine coordinates on $\mathcal Q_n$ and
$\Phi_r(p)\ge r^{-(n-1)}\sum_{s<n}|p(s)|$, so the sublevel sets of $\Phi_r$ are
bounded; $\Phi_r$ is convex and finite on $\mathcal Q_n$, hence continuous,
and a minimizer exists.  Let $p_*$ be any minimizer.  As noted before
Lemma~\ref{lem:vertexopt}, $\Psi_r(p_*)>0$ and $R_{p_*}(\sigma)\to0$, so $\sigma_*$ exists.  Moreover
$R_{p_*}(\sigma)-R_{p_*}(\sigma+1)=\sum_{s>\sigma}|p_*(s)|r^{-s}>0$, since
$p_*\ne0$ has finitely many zeros.  So $\Psi_r(p_*)-2R_{p_*}(\sigma)\ge0$ for
$\sigma\ge\sigma_*$, strictly for $\sigma>\sigma_*$, and (a) with $p=p_*$
gives the first claim.  By Proposition~\ref{prop:rowvalue}(b) and (c),
$V_r(n+1,2)=\inf_\sigma1/\eta_\sigma(n)$ and
$V_r(n+1,2)\le\beta_r(n)=1/\nu_1(n)$, while the positions
$\sigma\ge\sigma_*$ contribute terms at least $1/\nu_1(n)$.

(c) The prefix identity at $k=2$ reads
$V_r(n+1,2)=\min\{1/\nu_1(n),1/\nu_2(n)\}$.  If $\sigma_*>1$, then
$1/\nu_2(n)=1/\eta_1(n)$ is one of the terms in (b); if $\sigma_*=1$, then
$\nu_2(n)=\eta_1(n)\le\nu_1(n)$ by (b).  So by (b) the identity fails
exactly when some term with $2\le\sigma<\sigma_*$ is smaller than both.
\end{proof}

Part (a) needs no optimality: any $p\in\mathcal Q_n$ and any $\sigma_0$
with $2R_p(\sigma_0)\le\Psi_r(p)$ give $\eta_\sigma(n)\le\Phi_r(p)$ for all
$\sigma\ge\sigma_0$, since $R_p$ decreases.  The next proposition handles
the far positions this way.

\paragraph{Names of zero sets.}
In Propositions~\ref{prop:secondrowexact}, \ref{prop:secondrowhyp}
and~\ref{prop:thirdrowexact} each part fixes a root $r$ and a top-row degree $n$, and a named set of integers
records $n$ and the positions it must contain: $Y_n$ is a set of $n-1$
integers (a candidate for $\nu_1(n)$), $Y'_n$ is a set of $n$ integers
containing $1$, and, for a set $S$ of exempted positions, $Z_S(n)$ is a set
of $n+|S|-1$ integers containing $S$ (a candidate for $\mu_r(S,n+|S|)$),
written $Z_\sigma(n)$ and $Z_{\sigma,\sigma'}(n)$ for $S=\{\sigma\}$ and
$S=\{\sigma,\sigma'\}$.  No name is used at two different roots.  The names
fix only sizes and prescribed elements; every property of a named set is
proved where it is used.

\begin{proposition}[Second rows at the roots $\frac54$ and $\frac43$]\label{prop:secondrowexact}
Let $r=\frac54$ in (a)--(c), and let
\begin{align*}
  Z_2(15)&=\{1,2,5,8,12,17,22,28,35,44,53,65,78,94,115\},\\
  Y_{15}&=\{1,3,6,9,13,18,24,30,38,47,58,71,87,107\},\\
  Y_{13}&=\{1,3,6,10,15,21,27,36,45,57,72,91\}.
\end{align*}
\begin{enumerate}
\item[(a)] $V_r(16,2)=1/\Phi_r(Z_2(15))=1/\nu_3(14)=6.5694\ldots$, while
$\min_{i\le2}1/\nu_i(15)=\beta_r(15)=1/\Phi_r(Y_{15})=7.6492\ldots$.
\item[(b)] $V_r(17,3)\le1/\Phi_r(Z_2(15))<\beta_r(15)=\min_{i\le3}1/\nu_i(15)$.
So the prefix identity fails at $(\frac54,16,2)$ and at $(\frac54,17,3)$.
\item[(c)] $V_r(14,2)=\beta_r(13)=1/\Phi_r(Y_{13})=4.6297\ldots>1$ and
$\nu_2(13)<\nu_1(13)$.  So the prefix identity holds at $(\frac54,14,2)$
with its minimum at $i=1$: the second row is the top row of
$(x-\frac54)^{13}$.
\item[(d)] Let $r=\frac43$,
$Z_3(18)=\{1,2,3,6,8,11,15,18,23,28,34,41,48,57,66,78,91,109\}$ and
$Y_{18}=\{1,2,4,6,9,12,16,20,25,30,36,44,52,61,72,85,103\}$.  Then
$V_r(19,2)=1/\eta_3(18)=1/\Phi_r(Z_3(18))=1/\nu_4(16)=41.38\ldots$, while
$\eta_2(18)<\nu_1(18)$ and
$\min_{i\le2}1/\nu_i(18)=\beta_r(18)=1/\Phi_r(Y_{18})=42.34\ldots$.  So the
prefix identity fails at $(\frac43,19,2)$, and a worst exempted position
is $\sigma=3$, while $\sigma=2$ is not one.
\end{enumerate}
\end{proposition}

\begin{proof}
The proof is computer-assisted: every number below is an exact rational,
computed as described after Lemma~\ref{lem:vertexopt}, and every comparison
is verified by exact rational arithmetic in \texttt{tests/test\_rows.py}
(\texttt{test\_second\_row\_at\_five\_fourths\_fails} for (a) and (b),
\texttt{test\_second\_row\_at\_five\_fourths\_holds} for (c),
\texttt{test\_second\_row\_at\_four\_thirds} for (d)).  No floating
point is used, and the decimals in the statement are truncations of the
exact values, checked by exact comparison with the stated digits.  The
finite computation is this.  For a set $Z$ of integers, $\Phi_r(Z)$ is the
finite sum over $s\le\max Z$ plus the closed form of the series past
$\max Z$; the vertex test of Lemma~\ref{lem:vertexopt} for $Z$ and $S$
evaluates $G_y$, a finite sum plus such a closed form, and
$|e_y(y)|r^{-y}$ for each of the at most $|Z|$ elements $y\in Z\setminus
S$; the tail condition at $\sigma_0$ is the sign of
$2R_p(\sigma_0)-\Psi_r(p)$, evaluated in the same way; and for each
$\sigma<\sigma_0$ the test evaluates $\Phi_r$ on the listed candidate
sets through $\sigma$, at most $n+1$ for each $\sigma$, until one
satisfies the bound.  In (a), (c) and (d) $p=r_{Y_n}$, and $\sigma_0$ is the least
integer $\sigma\ge\max Y_n$ with $2R_p(\sigma)\le\Psi_r(p)$, found by
evaluating this sign at $\sigma=\max Y_n,\max Y_n+1,\ldots$; in all three
parts it is $\max Y_n$ itself.  Any $\sigma_0$ satisfying the tail condition
would do, and none of the proofs uses the least one.  The named sets are the
certificate: with them the statement can be rechecked by any exact
evaluation of $\Phi_r$ and $G_y$.

(a) Lemma~\ref{lem:vertexopt} with $S=\{2\}$ shows
$\eta_2(15)=\Phi_r(Z_2(15))$, with every inequality strict, so
$r_{Z_2(15)}$ is the unique minimizer, and with $S=\varnothing$ that
$\nu_1(15)=\Phi_r(Y_{15})$, which is less than $\Phi_r(Z_2(15))$.  The tail
condition $2R_p(\sigma_0)\le\Psi_r(p)$ holds for $p=r_{Y_{15}}$ at
$\sigma_0=\max Y_{15}=107$, so $\eta_\sigma(15)\le\Phi_r(Y_{15})$ for
$\sigma\ge107$ by the remark after Theorem~\ref{thm:secondrow}.  (Since
$r_{Y_{15}}$ is a minimizer for $\nu_1(15)$, Theorem~\ref{thm:secondrow}(b)
applies with $p_*=r_{Y_{15}}$, and its threshold $\sigma_*$, the least
$\sigma\ge1$ with the tail condition, is $25$, as recorded in
Section~\ref{sec:finiterows} and checked there by exact arithmetic.  The
argument here uses only $\sigma_0=107$, which is larger than necessary,
and not the value of $\sigma_*$.)  For every
$\sigma<107$ other than $2$, one of the following sets $Z\ni\sigma$ of $15$ integers has
$\Phi_r(Z)\le\Phi_r(Z_2(15))$:
$Y'_{15}=\{1,3,5,8,12,17,22,28,36,44,54,65,78,94,115\}$ if
$\sigma\in Y'_{15}$; $\mathrm{ins}(Y_{15},\sigma)$ if $\sigma\notin Y_{15}$;
$Y'_{15}$ with some element replaced by $\sigma$ (all $15$ replacements are
tried) if $\sigma\notin Y'_{15}$.
For each $\sigma$, \texttt{tests/test\_rows.py} generates these candidates
and checks that one of them satisfies the bound.  Hence
$\sup_\sigma\eta_\sigma(15)=\Phi_r(Z_2(15))$, and
Proposition~\ref{prop:rowvalue}(b) gives $V_r(16,2)=1/\Phi_r(Z_2(15))$.  Since
$[1,2]\subseteq Z_2(15)$,
$\eta_2(15)\le\nu_3(14)=\mu_r([1,2],16)\le\Phi_r(Z_2(15))=\eta_2(15)$.  Finally
$\nu_2(15)\le\Phi_r(Y'_{15})<\nu_1(15)$, as $1\in Y'_{15}$.

(b) Lemma~\ref{lem:rowmono} and (a) give the first inequality.  With (a),
the set $\{1,2,5,8,11,16,21,26,33,41,50,60,71,85,102,123\}\supseteq[1,2]$
gives $\nu_3(15)<\nu_1(15)$, so
$\min_{i\le3}1/\nu_i(15)=1/\nu_1(15)=\beta_r(15)$, which exceeds
$1/\Phi_r(Z_2(15))$ by (a).

(c) Lemma~\ref{lem:vertexopt} gives $\nu_1(13)=\Phi_r(Y_{13})<1$.  The
tail condition holds for $p=r_{Y_{13}}$ at $\sigma_0=\max Y_{13}=91$, and every
$\sigma<91$ lies in a set of $13$ integers with
$\Phi_r\le\Phi_r(Y_{13})$, found among the candidates of (a) built from
$Y_{13}$ and
$Y'_{13}=\{1,3,6,9,14,19,25,33,41,52,64,79,99\}$.  So
$\sup_\sigma\eta_\sigma(13)=\nu_1(13)$, and
Proposition~\ref{prop:rowvalue}(b) gives $V_r(14,2)=\beta_r(13)$.  Since
$1\in Y'_{13}$, $\nu_2(13)\le\Phi_r(Y'_{13})<\nu_1(13)$.

(d) As in (a): Lemma~\ref{lem:vertexopt} gives $\eta_3(18)=\Phi_r(Z_3(18))$,
with every inequality strict, so $r_{Z_3(18)}$ is the unique minimizer,
and $\nu_1(18)=\Phi_r(Y_{18})<\Phi_r(Z_3(18))$; the tail condition holds for
$p=r_{Y_{18}}$ at $\sigma_0=\max Y_{18}=103$; and every $\sigma<103$ other than $3$
lies in a set of $18$ integers with $\Phi_r\le\Phi_r(Z_3(18))$, found among
the candidates of (a) built from $Y_{18}$ and
$Y'_{18}=\{1,2,4,6,8,11,15,19,23,28,34,41,48,57,66,78,91,109\}$.  Since
$\{1,2\}\subseteq Y'_{18}$, $\eta_2(18)$ and $\nu_2(18)$ are at most
$\Phi_r(Y'_{18})<\nu_1(18)$, and $[1,3]\subseteq Z_3(18)$ gives
$\eta_3(18)=\nu_4(16)=\mu_r([1,3],19)$ as in (a).
\end{proof}

In (a) and (d) the second row equals a prefix value of a longer prefix and
a smaller degree than the prefix identity allows, $1/\nu_3(14)$ and
$1/\nu_4(16)$, although only one position is exempted.  The reason is that
the row is attained at a position $\sigma$ for which the unique minimizer
$r_Z$ of $\eta_\sigma(n)$ has all of $[1,\sigma]$ among its zeros, so that
$\eta_\sigma(n)\le\mu_r([1,\sigma],n+1)\le\Phi_r(Z)=\eta_\sigma(n)$ and the
row is $1/\mu_r([1,\sigma],n+1)=1/\nu_{\sigma+1}(n+1-\sigma)$.  Neither
Theorem~\ref{thm:onepolyrows} nor Theorem~\ref{thm:longdiag} reaches (c),
since both place the minimum at $i=k$.

\paragraph{The hypothesis $\nu_2(n)\ge\nu_1(n)$.}
In Proposition~\ref{prop:secondrowexact}(a) and~(d) the identity fails
with $\nu_2(n)<\nu_1(n)$.  When
$\nu_2(n)>\nu_1(n)$ the position $\sigma=2$ cannot break it:
$\eta_2(n)\le\nu_2(n)$.  Indeed, by Lemma~\ref{lem:vertexopt} some $Y$ with
$|Y|=n-1$ has $\Phi_r(Y)=\nu_1(n)$, and some $Z\ni1$ with $|Z|=n$ has
$\Phi_r(Z)=\nu_2(n)=\mu_r(\{1\},n+1)$.  If $1\in Y$, then $r_Y$ is
admissible for $\mu_r(\{1\},n+1)$ and $\nu_2(n)\le\nu_1(n)$.  So
$Y\subset[2,\infty)$ and $r_Y(2)\ge0$, while
$r_Z(2)=-\prod_{z\in Z\setminus\{1\}}(1-2/z)\le0$, and
Lemma~\ref{lem:crossing}(a) with $T=\varnothing$, $\sigma=2$, $L=n+1$ and
$\varphi=\nu_2(n)$ gives $\eta_2(n)=\mu_r(\{2\},n+1)\le\nu_2(n)$.  The same
argument covers every $\sigma$ at which $r_Y$ and $r_Z$ take opposite
signs, but not every $\sigma$, and the hypothesis does not suffice.

\begin{proposition}[A second-row failure with $\nu_2(n)>\nu_1(n)$]\label{prop:secondrowhyp}
Let $r=\frac{177}{167}$ and
\begin{align*}
  Y_{50}&=\{2,4,7,10,15,20,27,34,42,51,61,71,83,95,109,123,138,155,172,190,\\
  &\qquad210,230,252,275,299,324,350,378,407,437,469,502,537,574,613,654,696,\\
  &\qquad741,789,839,893,950,1011,1076,1148,1226,1314,1416,1543\},\\
  Y'_{50}&=\{1,4,6,10,15,20,26,33,41,50,59,70,81,93,107,121,136,152,169,187,\\
  &\qquad205,225,247,269,292,316,342,369,397,427,458,490,524,560,597,637,678,\\
  &\qquad721,767,815,866,920,978,1039,1105,1177,1256,1345,1448,1576\},\\
  Z_5(50)&=\{2,4,5,10,14,20,26,33,41,50,59,70,81,93,106,121,136,152,168,186,\\
  &\qquad205,225,246,269,292,316,342,369,397,427,458,490,524,560,597,637,678,\\
  &\qquad721,767,815,866,920,978,1039,1105,1177,1256,1345,1448,1576\}.
\end{align*}
Then $\nu_1(50)=\Phi_r(Y_{50})<\nu_2(50)=\Phi_r(Y'_{50})<
\eta_5(50)=\Phi_r(Z_5(50))$, the reciprocals being $4.2262\ldots$,
$4.2221\ldots$ and $4.2112\ldots$.  So
\[
  V_r(51,2)\le\frac1{\eta_5(50)}<\frac1{\nu_2(50)}=\min_{i\le2}\frac1{\nu_i(50)},
\]
and the prefix identity fails at $(\frac{177}{167},51,2)$ although
$\nu_2(50)>\nu_1(50)$.
\end{proposition}

\begin{proof}
The proof is computer-assisted, by exact rational arithmetic as in
Proposition~\ref{prop:secondrowexact}, in
\texttt{test\_second\_row\_hypothesis\_fails} of
\texttt{tests/test\_rows.py}; no floating point is used, and the decimals
are truncations of exact values.  Lemma~\ref{lem:vertexopt} with
$S=\varnothing$, $\{1\}$ and $\{5\}$ gives the three equalities, recalling
$\nu_2(50)=\mu_r(\{1\},51)$; for $Z_5(50)$ every inequality of the lemma is
strict, so $r_{Z_5(50)}$ is the unique minimizer for $\eta_5(50)$.  The two
inequalities between the values are exact comparisons.  Finally
Proposition~\ref{prop:rowvalue}(b) with $S=\{5\}$ gives
$V_r(51,2)\le1/\mu_r(\{5\},51)=1/\eta_5(50)$.
\end{proof}

The unique minimizer for $\eta_5(50)$ vanishes at $2$, $4$ and $5$ but not
at $1$ or $3$.  So the mechanism of Proposition~\ref{prop:secondrowexact}(a)
and~(d), a minimizer for $\eta_\sigma(n)$ vanishing on all of $[1,\sigma]$,
is absent here.  The margins are small: $\nu_2(50)$ exceeds $\nu_1(50)$ by less
than $0.1$ percent, and $\eta_5(50)$ exceeds $\nu_2(50)$ by less than
$0.3$ percent.  We have not determined $V_r(51,2)$.

Floating-point experiments, not used in any proof, suggest that such
failures need $r$ near $1$ and $n$ large.  On the rationals
$r=p/q\in(1,2)$ with $q\le10$ and $n\le14$, and on those in $[\frac54,2)$
with $q\le20$ for every $n$ up to the third consecutive one with
$\nu_2(n)<\nu_1(n)$, every $n$ with $\nu_2(n)\ge\nu_1(n)$ gave
$\max_{\sigma\ge2}\eta_\sigma(n)<0.92\,\nu_2(n)$.  At the root $r_n$ where
$\nu_2(n)=\nu_1(n)$, which decreases from $1.52\ldots$ at $n=4$ to
$1.05\ldots$ at $n=58$, the ratio $\max_{\sigma\ge2}\eta_\sigma(n)/\nu_1(n)$,
computed for $4\le n\le28$ and $n=30,34,\ldots,58$, rises from $0.84$ at
$n=4$ to $0.9998$ at $n=42$, $1.0017$ at $n=46$ and $1.0054$ at $n=58$,
with the maximum at $\sigma=3$ for $n\le12$ and at $\sigma=5$ from $n=13$
on; Proposition~\ref{prop:secondrowhyp} takes $r$ just below $r_{50}$.
We do not know whether a worst exempted position, when one exists, is
always small, whether the identity fails for every $1<r<2$ or only near
$1$, or which condition on $r$ and $n$ gives it at $k=2$.

\subsection{Every row is a finite minimum}\label{sec:finiterows}

Theorem~\ref{thm:secondrow} rests on one fact: a zero placed far enough out
does not increase $\Phi_r$.  For $k\ge3$ several exempted positions can
move out together, and a threshold for one far zero does not serve two.  At
$r=\frac54$ and $p_*=r_{Y_{15}}$ of Proposition~\ref{prop:secondrowexact}
the threshold of Theorem~\ref{thm:secondrow} is $\sigma_*=25$, while the
zeros $25$ and $26$ together raise $\Phi_r(p_*)$ by more than half
(exact rational arithmetic, \texttt{test\_far\_zeros} in
\texttt{tests/test\_rows.py}).  Lemma~\ref{lem:farzeros} gives a threshold
for up to $m$ far zeros at once.  It depends on the exempted positions below
it, and the positions that stay below the thresholds form a finite tree.  So
each row is the minimum of the row before it on its diagonal and finitely
many exempted-set terms (Theorem~\ref{thm:finiterows}).

For integers $m,s\ge1$, real $t>0$ and a real polynomial $q$ put
\[
  \phi_m(s,t)=\begin{cases}-s,&s\le t,\\(s-t)(s/t)^{m-1}-t,&s>t,\end{cases}
  \qquad
  \Delta_{q,m}(t)=\sum_{s\ge1}\phi_m(s,t)\,|q(s)|\,r^{-s}.
\]
Since $\phi_1(s,t)=-s+2(s-t)^+$, we have
$\Delta_{q,1}(\sigma)=2R_q(\sigma)-\Psi_r(q)$.

\begin{lemma}[Several far zeros]\label{lem:farzeros}
Let $r>1$, $m\ge1$, and let $q$ be a nonzero real polynomial.
\begin{enumerate}
\item[(a)] $\Delta_{q,m}(t)$ does not increase in $t>0$,
$\Delta_{q,m}\le\Delta_{q,m+1}$, and $\Delta_{q,m}(t)\to-\Psi_r(q)<0$ as
$t\to\infty$.  So there is a least integer $t_m(q)\ge1$ with
$\Delta_{q,m}(t_m(q))\le0$; moreover $\Delta_{q,m}(t)\le0$ for all real
$t\ge t_m(q)$, and $t_m(q)\le t_{m+1}(q)$.
\item[(b)] If $t>0$, $\Delta_{q,m}(t)\le0$, and $U\subset\N_{>0}$ is finite
with $|U|\le m$ and all elements at least $t$, then
$\Phi_r(q\,r_U)\le\Phi_r(q)$.
\end{enumerate}
For $m=1$, $t_1(p_*)$ is the threshold $\sigma_*$ of
Theorem~\ref{thm:secondrow}.
\end{lemma}

\begin{proof}
(a) Fix $s$.  For $t\ge s$, $\phi_m(s,t)=-s$.  For $0<t<s$, the product
$(s-t)(s/t)^{m-1}$ of a positive decreasing and a positive nonincreasing
function of $t$ decreases, so $\phi_m(s,t)$ decreases in $t$ there, and it
tends to $-s$ as $t\uparrow s$.  So $\phi_m(s,\cdot)$ is nonincreasing.  For
$s>t$, $(s/t)^{m-1}\le(s/t)^m$, so $\phi_m\le\phi_{m+1}$.  For fixed $t$,
$|\phi_m(s,t)|$ is at most a polynomial in $s$, so every series here
converges absolutely, and the first two claims follow term by term.  For
$t\ge1$, $|\phi_m(s,t)|\le s^m+s$, and $\sum_s(s^m+s)|q(s)|r^{-s}<\infty$, so
dominated convergence with $\phi_m(s,t)\to-s$ gives the limit.  Finally
$\Psi_r(q)>0$, since $q$ has finitely many zeros.

(b) Let $U\ne\varnothing$ and $u=\min U\ge t$, so $\Delta_{q,m}(u)\le0$ by
(a).  Every $v\in U$ has $|1-s/v|\le\max\{1,s/v\}\le\max\{1,s/u\}$, so
$|r_U(s)|\le|1-s/u|\max\{1,s/u\}^{m-1}$ for $s\ge1$.  Hence
$u(|r_U(s)|-1)\le\phi_m(s,u)$: for $s\le u$ the left side is at most
$(u-s)-u$, and for $s>u$ at most $(s-u)(s/u)^{m-1}-u$.  Multiplying by
$|q(s)|r^{-s}$ and summing gives
$u\bigl(\Phi_r(q\,r_U)-\Phi_r(q)\bigr)\le\Delta_{q,m}(u)\le0$.  For $m=1$,
$\Delta_{p_*,1}(\sigma)\le0$ is the condition $2R_{p_*}(\sigma)\le\Psi_r(p_*)$
defining $\sigma_*$.
\end{proof}

\begin{theorem}[Every row is a finite minimum]\label{thm:finiterows}
Let $r>1$ and $n,m\ge1$.
\begin{enumerate}
\item[(a)] For each finite $N\subset\N_{>0}$ with $|N|<m$ fix a minimizer
$q_N$ of $\mu_r(N,n+|N|)$ (Lemma~\ref{lem:vertexopt}) and put
$T(N)=t_{m-|N|}(q_N)$.  Let $\mathcal S$ be the set of
$S=\{\sigma_1<\cdots<\sigma_m\}\subset\N_{>0}$ with
$\sigma_{j+1}<T(\{\sigma_1,\ldots,\sigma_j\})$ for $0\le j<m$.  Then
$\mathcal S$ is finite, and
\[
  V_r(n+m,m+1)=\min\Bigl\{V_r(n+m-1,m),\
  \min_{S\in\mathcal S}\frac1{\mu_r(S,n+m)}\Bigr\},
\]
where the inner minimum is omitted if $\mathcal S=\varnothing$.
\item[(b)] Let $B>0$, and let $\mathcal T$ be a set of finite subsets of
$\N_{>0}$ of size less than $m$, with $\varnothing\in\mathcal T$.  To each
$N\in\mathcal T$ attach a real polynomial $q_N$ with $\deg q_N\le n+|N|-1$,
$q_N(0)=1$, $q_N|_N=0$ and $\Phi_r(q_N)\le B$, and an integer $T_N\ge1$ with
$\Delta_{q_N,m-|N|}(T_N)\le0$.  Suppose that $N\cup\{\sigma\}\in\mathcal T$
whenever $N\in\mathcal T$, $|N|\le m-2$ and $\max N<\sigma<T_N$
($\max\varnothing=0$), and that $\mu_r(N\cup\{\sigma\},n+m)\le B$ whenever
$N\in\mathcal T$, $|N|=m-1$ and $\max N<\sigma<T_N$.  Then
$\mu_r(S,n+m)\le B$ for every $S\subset\N_{>0}$ with $|S|=m$, and
$V_r(n+m,m+1)\ge1/B$.
\end{enumerate}
\end{theorem}

\begin{proof}
Both parts rest on one step.  Let $S=\{\sigma_1<\cdots<\sigma_m\}$ and
put $N_j=\{\sigma_1,\ldots,\sigma_j\}$ for $0\le j\le m$.  Fix $j<m$ and
let $q$ be real with $\deg q\le n+j-1$, $q(0)=1$, $q|_{N_j}=0$ and
$\Delta_{q,m-j}(\sigma_{j+1})\le0$.  The set
$U=S\setminus N_j$ has $m-j$ elements, all at least $\sigma_{j+1}$, so
Lemma~\ref{lem:farzeros}(b) gives $\Phi_r(q\,r_U)\le\Phi_r(q)$.  The
polynomial $q\,r_U$ has degree at most $n+m-1$, takes the value $1$ at $0$
and vanishes on $S$, so
\begin{equation}\label{eq:farstep}
  \mu_r(S,n+m)\le\Phi_r(q).
\end{equation}

(a) A set in $\mathcal S$ with given $\sigma_1,\ldots,\sigma_j$ has fewer
than $T(N_j)$ choices of $\sigma_{j+1}$, so $\mathcal S$ is finite by
induction on $j$.  Proposition~\ref{prop:rowvalue}(b) gives
$V_r(n+m,m+1)\le1/\mu_r(S,n+m)$ for each $S$, and Lemma~\ref{lem:rowmono}
gives $V_r(n+m,m+1)\le V_r(n+m-1,m)$.  Conversely let $|S|=m$ with
$S\notin\mathcal S$.  Then some $j<m$ has $\sigma_{j+1}\ge T(N_j)$, so
$\Delta_{q_{N_j},m-j}(\sigma_{j+1})\le0$ by Lemma~\ref{lem:farzeros}(a),
and~\eqref{eq:farstep} gives $\mu_r(S,n+m)\le\mu_r(N_j,n+j)$.  By
Proposition~\ref{prop:rowvalue}(b), $1/\mu_r(N_j,n+j)\ge V_r(n+j,j+1)$ (for
$j=0$ both sides are $\beta_r(n)$), and by Lemma~\ref{lem:rowmono},
$V_r(n+j,j+1)\ge V_r(n+m-1,m)$ since $j\le m-1$.  So every $S$ has
$1/\mu_r(S,n+m)$ at least the right side, and
Proposition~\ref{prop:rowvalue}(b) gives the reverse inequality.

(b) Let $|S|=m$, and let $j<m$ be least with $\sigma_{j+1}\ge T_{N_j}$, if
any; then $N_0,\ldots,N_j\in\mathcal T$ by induction, since
$N_0=\varnothing\in\mathcal T$, and $N_i\in\mathcal T$,
$\sigma_{i+1}<T_{N_i}$, $i\le m-2$ give $N_{i+1}\in\mathcal T$.  If there
is none, then likewise $N_{m-1}\in\mathcal T$ and $\sigma_m<T_{N_{m-1}}$,
so $\mu_r(S,n+m)\le B$ by the leaf hypothesis.  Otherwise
$\Delta_{q_{N_j},m-j}(\sigma_{j+1})\le\Delta_{q_{N_j},m-j}(T_{N_j})\le0$ by
Lemma~\ref{lem:farzeros}(a), and~\eqref{eq:farstep} gives
$\mu_r(S,n+m)\le\Phi_r(q_{N_j})\le B$.  Proposition~\ref{prop:rowvalue}(b)
gives $V_r(n+m,m+1)\ge1/B$.
\end{proof}

For $m=1$ the set $\mathcal S$ is $\{\{\sigma\}:\sigma<\sigma_*\}$ and (a) is
Theorem~\ref{thm:secondrow}(b).  Applying (a) to $V_r(n+m-1,m)$, and so on
down to $V_r(n,1)=\beta_r(n)$, writes every row $V_r(n+m,m+1)$ as the
minimum of finitely many values $1/\mu_r(S,n+|S|)$, $|S|\le m$, each equal
to $1/\Phi_r(Z)$ for a set $Z\supseteq S$ of $n+|S|-1$ integers
(Lemma~\ref{lem:vertexopt}).  The sets in $\mathcal S$ depend on the
minimizers $q_N$, and the theorem does not say how to find them; part (b)
needs only upper bounds.  If $V_r(n+m,m+1)<V_r(n+m-1,m)$, the infimum
defining $V_r(n+m,m+1)$ is attained, at some $S\in\mathcal S$.  For rational
$r$, $q=r_Z$ and an integer $t\ge1$, $\Delta_{q,m}(t)$ is a rational number
computable exactly
as $\Phi_r(Z)$ is, since past $\max(Z\cup\{t\})$ its summand is a polynomial
times $r^{-s}$.  So for given zero sets every hypothesis of (b), with
$\mu_r(N\cup\{\sigma\},n+m)\le B$ checked through some $\Phi_r(Z)\le B$,
$Z\supseteq N\cup\{\sigma\}$, $|Z|\le n+m-1$, is a finite exact check.

\begin{proposition}[Third rows at the roots $\frac43$ and $\frac32$]\label{prop:thirdrowexact}
Let $r=\frac43$ in (a) and (b).
\begin{enumerate}
\item[(a)] Let
\begin{align*}
  Y_{12}&=\{1,3,6,9,13,17,23,30,39,50,64\},\\
  Z_2(12)&=\{1,2,5,8,12,16,21,28,35,44,56,71\},\\
  Z_{2,3}(12)&=\{2,3,4,7,10,15,19,25,32,40,49,61,77\}.
\end{align*}
Then $V_r(14,3)=1/\Phi_r(Z_{2,3}(12))=7.7168\ldots$,
$V_r(13,2)=1/\Phi_r(Z_2(12))=1/\nu_3(11)=8.5482\ldots$ and
$\beta_r(12)=1/\Phi_r(Y_{12})=8.8501\ldots$, and
$\min_{i\le k}1/\nu_i(12)=\beta_r(12)$ for $k=2,3,4$.  So
$V_r(14,3)<V_r(13,2)<\beta_r(12)$: the prefix identity fails at
$(\frac43,13,2)$ and at $(\frac43,14,3)$, and the third row lies strictly
below the second on this diagonal.  By Lemma~\ref{lem:rowmono} it also
fails at $(\frac43,15,4)$.  The pair $S=\{2,3\}$ is the only
exempted set with $1/\mu_r(S,14)=V_r(14,3)$.
\item[(b)] Let
\begin{align*}
  Y_{14}&=\{1,3,5,8,11,15,20,26,32,40,50,62,77\},\\
  Z_{2,3}(14)&=\{1,2,3,6,9,13,17,22,27,34,41,50,60,73,89\}.
\end{align*}
Then $V_r(15,2)=\beta_r(14)=1/\Phi_r(Y_{14})=15.0390\ldots=\min_{i\le4}1/\nu_i(14)$
and $V_r(16,3)=1/\Phi_r(Z_{2,3}(14))=1/\nu_4(13)=14.6268\ldots$.  So the
prefix identity holds at $(\frac43,15,2)$, with its minimum at $i=1$, and
fails at $(\frac43,16,3)$ and, by Lemma~\ref{lem:rowmono}, at
$(\frac43,17,4)$.  The pair $S=\{2,3\}$ is the only exempted set
with $1/\mu_r(S,16)=V_r(16,3)$.
\item[(c)] Let $r=\frac32$ and $Y_8=\{1,3,6,9,14,20,28\}$.  Then
$V_r(10,3)=V_r(9,2)=\beta_r(8)=1/\Phi_r(Y_8)=8.0259\ldots=\min_{i\le3}1/\nu_i(8)$,
and $\nu_2(8),\nu_3(8)<\nu_1(8)$.  So the prefix identity holds at
$(\frac32,10,3)$ with its minimum at $i=1$.
\end{enumerate}
\end{proposition}

\begin{proof}
The proof is computer-assisted, by exact rational arithmetic as in
Proposition~\ref{prop:secondrowexact}, in \texttt{tests/test\_rows.py}
(\texttt{test\_third\_row\_below\_second} for (a),
\texttt{test\_third\_row\_fails\_alone} for (b),
\texttt{test\_third\_row\_is\_top\_row} for (c)); no floating point is
used, and the decimals are truncations of exact values.  Each lower bound for a row comes from
Theorem~\ref{thm:finiterows}(b) with $q_\varnothing=r_Y$ for the set $Y$
named in that part, $q_N=r_{K_N}$ for sets $K_N\supseteq N$ of at most
$n+|N|-1$ integers, and $T_N$ the least integer with
$\Delta_{q_N,m-|N|}(T_N)\le0$; a leaf $N\cup\{\sigma\}$ is checked through a
set $Z\supseteq N\cup\{\sigma\}$ of at most $n+m-1$ integers with
$\Phi_r(Z)\le B$.  The test builds the tree from the root
$K_\varnothing=Y$.  For a child $N\cup\{\sigma\}$ of a node $N$ with kept
set $K_N$ it forms these candidates: $K_N$ itself if $\sigma\in K_N$;
otherwise $K_N$ with one of the three elements of $K_N\setminus N$ nearest
to $\sigma$ replaced by $\sigma$, and $K_N\cup\{\sigma\}$; and in either
case the fixed set of $n+|N|$ integers of that part, with each position of
$N\cup\{\sigma\}$ that it misses, taken in increasing order, put in place of
one of the two elements outside $N\cup\{\sigma\}$ nearest to it, in every
combination.  The fixed sets are
\begin{align*}
  \text{(a)}\quad&\{1,3,5,8,12,16,21,28,35,44,56,71\},\
  \{1,3,5,8,11,15,20,26,32,40,50,62,77\},\\
  \text{(b)}\quad&\{1,3,5,7,10,14,18,24,30,37,45,55,67,83\},\
  \{1,3,4,7,10,13,17,22,28,34,42,50,61,73,90\},\\
  \text{(c)}\quad&\{1,3,5,8,12,17,24,33\},\ \{1,3,5,8,11,15,21,28,37\},
\end{align*}
of $n$ and $n+1$ integers, used for children of size one and two.  Among
the candidates with at most $n+|N|$ elements the test keeps one with the
least $\Phi_r$ (on ties, the first formed, and nearest elements are taken
toward the smaller one), unless the child is given a named set.  Each
threshold $T_N$ is found by evaluating the sign of $\Delta_{q_N,m-|N|}$ at
integers (it does not increase, Lemma~\ref{lem:farzeros}(a)), and the test
checks $\Delta_{q_N,m-|N|}(T_N)\le0$ exactly.  For the third rows the trees
have $32$, $35$ and $15$ nodes of size one and $126$, $161$ and $33$ leaves
in (a), (b) and (c); at each node and leaf the test checks $\Phi_r\le B$,
or $\Phi_r<B$ where stated, for the set it keeps.  Theorem~\ref{thm:finiterows}(b)
needs only that each node and each leaf have some set of the allowed size
through it with $\Phi_r\le B$ (or $<B$); which one is kept, and so the
tie-breaking above, changes the thresholds $T_N$ and the tree sizes but not
the validity of the certificate.

(a) Lemma~\ref{lem:vertexopt} with $S=\{2,3\}$, $\{2\}$ and $\varnothing$
gives $\mu_r(\{2,3\},14)=\Phi_r(Z_{2,3}(12))$, $\eta_2(12)=\Phi_r(Z_2(12))$ and
$\nu_1(12)=\Phi_r(Y_{12})$, and $\Phi_r(Y_{12})<\Phi_r(Z_2(12))<\Phi_r(Z_{2,3}(12))$.
For $Z_{2,3}(12)$ every inequality of the lemma is strict, so
$r_{Z_{2,3}(12)}$ is the unique minimizer for $\mu_r(\{2,3\},14)$.
With $n=12$, $m=2$, $B=\Phi_r(Z_{2,3}(12))$, $Y=Y_{12}$ (so
$T_\varnothing=33$) and the pair $\{2,3\}$ given $Z_{2,3}(12)$, the test checks
the hypotheses of Theorem~\ref{thm:finiterows}(b), with $\Phi_r<B$ for every
other set used.  So $\mu_r(S,14)\le B$ for all $|S|=2$, with equality only
at $S=\{2,3\}$: a pair outside the tree has $\mu_r(S,14)\le\Phi_r(q_N)<B$
by~\eqref{eq:farstep}.  Proposition~\ref{prop:rowvalue}(b) gives
$V_r(14,3)=1/B$.  With $m=1$, $B=\Phi_r(Z_2(12))$ and the position $2$ given
$Z_2(12)$, the same check gives $V_r(13,2)\ge1/\Phi_r(Z_2(12))$, and $S=\{2\}$ gives
equality.  Since $[1,2]\subseteq Z_2(12)$,
$\eta_2(12)\le\nu_3(11)=\mu_r([1,2],13)\le\Phi_r(Z_2(12))=\eta_2(12)$.  Finally
the sets $\{1,3,5,8,12,16,21,28,35,44,56,71\}\ni1$,
$\{1,2,5,7,11,15,20,25,32,40,50,61,77\}\supseteq[1,2]$ and
$\{1,2,3,7,10,14,18,23,29,37,45,55,67,83\}\supseteq[1,3]$ give
$\nu_i(12)<\nu_1(12)$ for $i=2,3,4$, so
$\min_{i\le4}1/\nu_i(12)=1/\nu_1(12)$, and Lemma~\ref{lem:rowmono} gives
$V_r(15,4)\le V_r(14,3)$.

(b) As in (a): Lemma~\ref{lem:vertexopt} gives
$\mu_r(\{2,3\},16)=\Phi_r(Z_{2,3}(14))$, with every inequality strict, so
$r_{Z_{2,3}(14)}$ is the unique minimizer, and $\nu_1(14)=\Phi_r(Y_{14})$, which is
smaller.  Theorem~\ref{thm:finiterows}(b) with $n=14$, $m=2$,
$B=\Phi_r(Z_{2,3}(14))$ and $Y=Y_{14}$ ($T_\varnothing=36$), with $\Phi_r<B$
for every set used except $Z_{2,3}(14)$ at $\{2,3\}$, gives $V_r(16,3)=1/B$,
attained only at $\{2,3\}$.  With $m=1$ and $B=\Phi_r(Y_{14})$ it gives
$V_r(15,2)\ge\beta_r(14)$, and Proposition~\ref{prop:rowvalue}(c) gives
equality.  Since $[1,3]\subseteq Z_{2,3}(14)$,
$\mu_r(\{2,3\},16)\le\mu_r([1,3],16)=\nu_4(13)\le\Phi_r(Z_{2,3}(14))$, so all
three are equal.  The sets $\{1,3,5,7,10,14,18,24,30,37,45,55,67,83\}\ni1$,
$\{1,2,4,7,10,13,17,22,28,34,42,50,61,73,90\}\supseteq[1,2]$ and
$\{1,2,3,6,9,12,16,21,26,32,39,46,55,66,79,96\}\supseteq[1,3]$ give
$\nu_i(14)<\nu_1(14)$ for $i=2,3,4$, and Lemma~\ref{lem:rowmono} gives
$V_r(17,4)\le V_r(16,3)$.

(c) Lemma~\ref{lem:vertexopt} gives $\nu_1(8)=\Phi_r(Y_8)$.
Theorem~\ref{thm:finiterows}(b) with $n=8$, $m=2$, $B=\Phi_r(Y_8)$, $Y=Y_8$,
and the pair $\{2,3\}$ given $\{2,3,4,7,10,15,20,27,37\}$, gives
$V_r(10,3)\ge\beta_r(8)$.  Lemma~\ref{lem:rowmono} and
Proposition~\ref{prop:rowvalue}(c) give
$V_r(10,3)\le V_r(9,2)\le\beta_r(8)$.  The sets
$\{1,3,5,8,12,17,24,33\}\ni1$ and $\{1,2,5,7,11,15,21,28,37\}\supseteq[1,2]$
give $\nu_2(8),\nu_3(8)<\nu_1(8)$.
\end{proof}

In (b) the prefix identity fails in the third row although it holds in the
second, so this failure does not come from the second row through
Lemma~\ref{lem:rowmono}.  In (a) the third row is strictly below the
second.  In (b) the zero set $Z_{2,3}(14)$ of the unique minimizer for
$\mu_r(\{2,3\},16)$ contains $[1,3]$, as the zero sets of the unique
minimizers for $\eta_\sigma(n)$ in Proposition~\ref{prop:secondrowexact}(a)
and~(d) contain $[1,\sigma]$; in (a) the zero set $Z_{2,3}(12)$ of the
unique minimizer for $\mu_r(\{2,3\},14)$ omits $1$.  The identities at
$(\frac43,15,2)$ and $(\frac32,10,3)$ are outside
Theorems~\ref{thm:onepolyrows} and~\ref{thm:longdiag}, which give
$\nu_1(n)\le\nu_k(n)$.

\paragraph{Data availability.}
No datasets were generated or analyzed in this study.  The exact
computations behind Propositions~\ref{prop:secondrowexact},
\ref{prop:secondrowhyp} and~\ref{prop:thirdrowexact} are in
\texttt{tests/test\_rows.py} of the repository
\href{https://github.com/bangyen/coefficient-mass}{bangyen/coefficient-mass};
they need Python~3.12 or later.

\paragraph{Computer assistance.}
Propositions~\ref{prop:secondrowexact}, \ref{prop:secondrowhyp}
and~\ref{prop:thirdrowexact} are the only computer-assisted results; every other result is proved by hand.  Their
proofs reduce each claim, through Lemma~\ref{lem:vertexopt},
Theorem~\ref{thm:secondrow}(a) and Theorem~\ref{thm:finiterows}(b), to
finitely many comparisons of explicit rational numbers: the vertex test of
Lemma~\ref{lem:vertexopt} for each named zero set, the sign of
$\Delta_{q,m}$ at each threshold, the bound $\Phi_r(Z)\le B$ (or $<B$) for
a zero set through each exempted position or set below the thresholds, and
the comparisons with the stated decimals.  The tests
\texttt{test\_second\_row\_*} (including
\texttt{test\_second\_row\_hypothesis\_fails}) and \texttt{test\_third\_row\_*} in
\texttt{tests/test\_rows.py} carry out these comparisons in exact rational
arithmetic (Python's \texttt{fractions} module), evaluating $\Phi_r$, $G_y$
and $\Delta_{q,m}$ through the closed form described after
Lemma~\ref{lem:vertexopt}, with no floating point; these computations are
steps of the proofs.  The zero sets were found by search, and only their
verification is part of the proofs.  The value $\sigma_*=25$ and the
comparison for the zeros $25$ and $26$ quoted in
Section~\ref{sec:finiterows} are computed in the same way
(\texttt{test\_far\_zeros}) and are illustrations, not used in any proof.
The floating-point experiments reported after
Proposition~\ref{prop:secondrowhyp} are not part of any proof or test; they
only guided the search for its zero sets.
The remaining tests are checks, not proof steps.
\texttt{test\_prefix\_identity\_fails\_below\_two} builds from
$Z_2(15)$, by complementary slackness, an explicit monic multiple of
$(x-\frac54)^{16}$ of degree $300$ with rational coefficients, verifies the
divisibility by $16$ exact synthetic divisions, and checks that its second
largest nonleading coefficient magnitude is at most $6.5695$, an
independent primal check of the upper bound in
Proposition~\ref{prop:secondrowexact}(a).  Other tests compare
Corollary~\ref{cor:lastrowr} and Theorem~\ref{thm:allrowstwo} with
exhaustive minima over zero sets in a window for small $L$, and test the
closed form for $\Phi_r$ and the vertex test on known cases.  Each check
is paired with a control that it must reject.  The whole file runs in a
few seconds.

\paragraph{Use of AI tools.}
Large language model assistants (Claude, Anthropic) were used in drafting
and revising the text, the proofs and the verification code. The author has checked the mathematics and is responsible
for the content.

\paragraph{Funding and competing interests.}
The author received no specific funding for this work and has no competing
interests to declare.

\begin{thebibliography}{9}

\bibitem{bbbp-multiple}
F.~Beaucoup, P.~Borwein, D.~W. Boyd and C.~Pinner,
\newblock Multiple roots of $[-1,1]$ power series,
\newblock \emph{Journal of the London Mathematical Society} (2) 57 (1998), 135--147,
\newblock \href{https://doi.org/10.1112/S0024610798005857}
{doi:10.1112/S0024610798005857}.

\bibitem{dj}
A.~Dubickas and J.~Jankauskas,
\newblock On the reduced height of a polynomial,
\newblock \emph{Publicationes Mathematicae Debrecen} 71 (2007), 325--348,
\newblock \href{https://doi.org/10.5486/PMD.2007.3728}
{doi:10.5486/PMD.2007.3728}.

\bibitem{coefficient-mass}
B.~Pham,
\newblock Displaced-zero certificates and the coefficient mass of polynomial
multiples,
\newblock preprint, 2026, \url{https://github.com/bangyen/coefficient-mass}.

\bibitem{schinzel}
A.~Schinzel,
\newblock On the reduced length of a polynomial with real coefficients,
\newblock \emph{Functiones et Approximatio Commentarii Mathematici} 35
(2006), 271--306,
\newblock \href{https://doi.org/10.7169/facm/1229442629}
{doi:10.7169/facm/1229442629}.

\end{thebibliography}

\end{document}
